\documentclass[11pt,reqno]{amsart}
\usepackage[T1]{fontenc}
\usepackage{lmodern,amsmath,amssymb,mathtools,microtype,booktabs,array,longtable}
\usepackage[margin=1in]{geometry}
\usepackage{graphicx}
\usepackage{xr-hyper}
\usepackage[hidelinks]{hyperref}
\ifdefined\pdfinfoomitdate\pdfinfoomitdate=1\fi
\ifdefined\pdftrailerid\pdftrailerid{}\fi
\ifdefined\pdfsuppressptexinfo\pdfsuppressptexinfo=15\fi
\numberwithin{equation}{section}
\newtheorem{theorem}{Theorem}[section]
\newtheorem{proposition}[theorem]{Proposition}
\newtheorem{lemma}[theorem]{Lemma}
\newtheorem{corollary}[theorem]{Corollary}
\theoremstyle{definition}\newtheorem{definition}[theorem]{Definition}
\newtheorem{example}[theorem]{Example}
\theoremstyle{remark}\newtheorem{remark}[theorem]{Remark}
\newcommand{\E}{\mathbb E}\newcommand{\Prb}{\mathbb P}
\newcommand{\R}{\mathbb R}\newcommand{\one}{\mathbf1}
\newcommand{\Q}{\mathcal Q}\newcommand{\TV}{\mathrm{TV}}
\newcommand{\cp}{\mathrm{cp}}\newcommand{\ev}{\mathrm{ev}}\newcommand{\st}{\mathrm{st}}
\DeclareMathOperator{\tr}{tr}\DeclareMathOperator{\dist}{dist}
\DeclareMathOperator{\law}{Law}\DeclareMathOperator{\rank}{rank}
\DeclareMathOperator{\supp}{supp}
\emergencystretch=2em

\externaldocument[][nocite]{supplement-refs}[General_Theta_Foundations_I_R71_Supplement.pdf]
\hypersetup{pdftitle={General Theta Foundations I: Acquired Geometry and Causal Resource Transfer},pdfauthor={Qian Qi}}
\title[Acquired geometry and causal resource transfer]{General Theta Foundations I:\\Acquired Geometry and Causal Resource Transfer}
\author{Qian Qi}\date{October 2026}
\subjclass[2020]{Primary 60G40, 62C20; Secondary 11J70, 68Q45, 81P15, 90C22, 94A29}
\keywords{causal experiment, exact memory, Diophantine approximation, stochastic realization, boundary optimality, quantum instrument, acquired geometry}
\begin{document}
\begin{abstract}
We study how the arithmetic of a fixed phase controls exact finite classical memory in a prepared causal experiment. For a damped single-frequency probability response, we prove a phase-uniform two-sided comparison between its least stochastic memory and an explicit rational-approximation envelope. The lower applies to arbitrary stochastic matrices and bounded decoders; the upper constructs a counted pure-start circulant observer. For every irrational phase of irrationality exponent $\mu$, the lower and upper logarithmic memory-growth exponents are exactly $1/(\mu+1)$ and $(\mu-1)/(2\mu-1)$, with limits $0$ and $1/2$ when $\mu$ is infinite. A uniform cubic-root law holds precisely for badly approximable phases. Every other irrational phase has normalized cubic-root memory with lower limit zero and upper limit infinity, even when its logarithmic exponent is $1/3$. Finite spectral mixtures admit the corresponding envelope law, and the inherited exact leakage increment transfers the single-phase conclusions without dependence on absolute leakage. These are exact optimal-observer counts for one fixed squared-loss embedding at its original least expected cost, not approximate simulation or an identification of physical dimension with memory. The proofs retain the original causal transfer, stopping, instrument-realization and spectral results with their separate hypotheses. Classical stochastic spectral geometry and continued-fraction identities supply the ingredients; the new results quantify their full fixed-phase memory consequences.
\end{abstract}
\maketitle
\section{Introduction and current theorem structure}\label{sec:current-overview}\label{sec:introduction}
A finite-dimensional prepared quantum system can have an optimal observable probability response whose exact classical implementation needs arbitrarily large memory. The required size depends not only on noise magnitude but on rational approximation of its phase. We distinguish the whole memory curve, its logarithmic growth exponents and a uniform power law. These are different mathematical questions even at one fixed phase.

The new main theorem, Theorem~\ref{thm:r71-envelope}, identifies an explicit arithmetic envelope comparable to the least exact stochastic memory with universal constants. It has no badly-approximable hypothesis. Theorem~\ref{thm:r71-exponents} determines both logarithmic exponents from the irrationality exponent of any fixed phase. Theorem~\ref{thm:r71-cubic-rigidity} proves that a uniform cubic-root law holds exactly for badly approximable phases, and proves unbounded fluctuations after cubic-root normalization at every other irrational phase. Corollary~\ref{cor:r71-mixtures} carries the envelope to finite spectral mixtures and leakage. All new statements have full proofs below.

The response-to-spectrum equivalence, stochastic sector inequality and its equality structure, rational polygon construction and exact leakage increment are inherited from the preceding complete manuscript. They are included after the new theorem to expose its complete positive-matrix foundation. Classical stochastic eigenvalue regions and positive-realization methods are not being renamed as new inventions \cite{DmitrievDynkin46,JohnsonPaparella17,KirklandSmigoc22,MonrasWinter14}. The continued-fraction identity used here is classical and is proved explicitly alongside its credit \cite{HoneVarona19}. Noise-dependent lower bounds for classical realizations of quantum processes, including a worst-angle square-root law, precede this paper \cite{Fanizza24}; the present conclusions fix an arbitrary phase and determine its entire approximation-sensitive envelope and two exact growth exponents.

The operational interpretation is unchanged. A real three-outcome audit with one fixed squared loss and two fixed decoder values forces the target probability response at the original least expected-cost boundary. No extra observation of the protected system, free initial mixture, uncharged clock or freely retained random seed is used. The new theorem does not infer a classical exact-memory rate merely from the physical Hilbert dimension or physical mixing speed.

The full current supplement preserves the stopping-rigidity, robust-tie, instrument-realization, finite-model witness and performed-law theorems and the original reset, rank, spectral, geometric, predictive and effective developments with their individual assumptions. The main article also retains the actual-history lower and common-program upper proof. The source retention map separates this paired current mathematics from all unchanged historical manuscripts and independent realization lines. A new phase theorem neither silently deletes those results nor makes every independent companion a premise of its proof.

\section{Prepared kernels, acquired task laws and observers}\label{sec:data}
All hidden-state, action and report spaces are standard Borel. A fixed unknown index $\theta\in\Theta$ specifies an actual preparation and positive normalized causal kernels. Positivity means nonnegative measures, not an everywhere strictly positive density. Dependence on the entire physical past may be included in the hidden state. A legal policy sees past visible reports and its admitted randomization, but not $\theta$ or a future audit. Iteration of these kernels defines its finite-prefix laws; the consistent countable path law is their kernel extension.

Each call has a cost $c_i\ge c_*>0$. The decision to make the next call is predictable. In the geometric route, the declared productive-call flag is also predictable; in the primitive-instrument route, productive reports are specified outcomes and can be random. Commitment is permitted only after a productive report, and occurs exactly once almost surely in each model. At that time a specified future audit $Z\in H$ is scored, where $H\subset\R^k$ is nonempty, compact and convex. The prediction is committed before that audit is observed. A prospective joint report--audit kernel is defined at every productive call, including rejected calls. The stopping coin is independent of the audit conditional on the entire current visible report and the past. Write $L=\operatorname{diam}(H)^2$; the case $L=0$ is trivial whenever a legal rule exists.

The core cost $C_\tau$ includes initial preparation charged to the core, failed calls and all productive calls through commitment. A separately executed final audit must be added to the raw budget. This convention is fixed before taking a risk infimum. Physical time, stored program constants and temporary computation are not identified with call count.

\subsection*{Three resource interfaces}
A checkpoint procedure may use its whole precommit history, but must output one element of a fixed dictionary of at most $M$ values. This counts its output width only; it is not a claim of $\log_2 M$ total working bits.

An autonomous observer has an alphabet $S$, $|S|\le M$, and one fixed parameter-independent program. At an intercall boundary its only persistent observer memory is its label $s\in S$. Its next command distribution is a fixed function of $s$. A physical port name, elapsed time, subprogram index, random seed or former report is not an extra free input. During a call its current command, current report and fresh coins are available. Whatever persists to another call must be encoded in $S$. A fixed decoder $r:S\to H$ reads only the committed label.

In the \emph{event-stopped} interface a post-report kernel emits $(s',e)\in S\times\{0,1\}$. An event $e=1$ ends the experiment and retains $s'$ for the decoder; no subsequent command is executed. This public commitment event is not a stored bit used in continuing computation. The same label may be used live on one path and terminal on another. In the \emph{state-stopped} interface a fixed map $h_*:S\to\{0,1\}$ determines halt status. The initial label is live, and type-one labels are globally absorbing terminal labels with no executed command rows. Entry into a terminal label is allowed only after a productive report. Thus live and terminal types are disjoint in this interface, and only in this interface by definition.

Write $R_i^\Theta(N,M)$, $i\in\{\cp,\ev,\st\}$, for the infimum of worst-case expected squared loss over programs with robust almost-sure commitment and $\sup_\theta\E_\theta C_\tau\le N$. An empty legal class has value $+\infty$. A program chosen once at random is not a free persistent seed. The ideal measurable class may have fixed real constants; effective implementation is stated separately in Section~\ref{sec:composition}.

\subsection*{A physically acquired comparison law}
Execute a declared prepared comparison experiment and obtain its joint report--audit probability $K_0(dx,dz)$ on $X\times H$. Its report marginal is $\nu$. The law is genuinely preparable and its specification has its own resource account. It need not be a member of $\Theta$. In particular no finite set of calibration data is silently promoted to a uniform certificate. Set
\begin{equation}\label{eq:derived}
 f(x)=\int zK_0(dz\mid x),\quad
 w(x)=\int\|z-f(x)\|^2K_0(dz\mid x),\quad
 d_r(x)=w(x)+\|f(x)-r\|^2.
\end{equation}
Disintegration gives Borel versions with $f\in H$; null sets may be completed arbitrarily. The profile is
\begin{equation}\label{eq:profile}
 \Q_m(b)=\min_{\substack{C\subset H,\ 1\le|C|\le m\\
                  0\le h\le b,\ \int h\,d\nu=1}}
        \int [w+\dist(f,C)^2]h\,d\nu,
 \qquad b\ge1,\ m\ge1.
\end{equation}
We use $+\infty$ for an infeasible cap or width. The quantities $f,w$ are task coordinates of acquired predictive information, not a complete substitute for the predictive quotient. Their derivation and the quotient construction in Appendix~\ref{sec:quotient} precede any use of local dimension.

\subsection*{Lower certificates on actual histories}\label{sec:lower-cert}
Fix \emph{one} actual model $\theta_*\in\Theta$, independent of the competing program. At each actual reached productive pre-call history $h$ and legal command $a$, write $P_{h,a}(dy,dz)$ for its conditional full current report--audit law. Let $P^0_{h,a}$ be a measurable comparison row on the same report space, with a Borel mark $x=\pi_{h,a}(y)$ of law $\nu$ and
\begin{equation}\label{eq:full-moments}
 \E^0[Z\mid y]=f(\pi_{h,a}y),\qquad
 \E^0[\|Z\|^2\mid y]=w(\pi_{h,a}y)+\|f(\pi_{h,a}y)\|^2.
\end{equation}
These are conditional identities given the \emph{full} report, not just its mark. Versions and inequalities are required on actual reachable histories, not only almost everywhere under an unrelated comparison path law.

Fix $0<l_-\le1\le u_-<\infty$. For all Borel $q:Y\to[0,1]$ and all finite-image Borel $r:Y\to H$, suppose predictable nonnegative defects $a_-(h,a),e_-(h,a)$ obey
\begin{align}
 \int q\,d(P_{h,a})_Y
 &\le u_-\int q\,d(P^0_{h,a})_Y+a_-(h,a),\label{eq:lower-mass}\\
 \int q(y)\|z-r(y)\|^2P_{h,a}(dy,dz)
 &\ge l_-\int q(y)\|z-r(y)\|^2P^0_{h,a}(dy,dz)-e_-(h,a).
 \label{eq:lower-task}
\end{align}
Independent current randomization can be adjoined to both rows: integrating inequalities for each fixed random value proves the same inequalities with the same defects. All command-conditioned lower inequalities hold before averaging the chosen command.

Let $A_\tau$ count productive calls. At this same actual model require, for every legal competing procedure,
\begin{equation}\label{eq:effort}
 \E_* C_\tau\ge\kappa_-\E_* A_\tau,\qquad\kappa_->0.
\end{equation}
At budget $N$ suppose, uniformly over all robustly feasible competitors,
\begin{equation}\label{eq:occupation}
 \E_*\sum_{j\le\tau:\ j\ {\rm productive}}a_-(H_{j-1},a_j)\le A_N,
 \qquad
 \E_*\sum_{j\le\tau:\ j\ {\rm productive}}e_-(H_{j-1},a_j)\le E_N.
\end{equation}
These are actual occupation certificates, not properties merely of the proposed upper program. Uniform row defects give $A_N\le(N/\kappa_-)\bar a_-$ and $E_N\le(N/\kappa_-)\bar e_-$. Lemma~\ref{lem:drift} supplies a nonuniform alternative.

\subsection*{One common upper construction}\label{sec:upper-cert}
Independently, let $\mathcal G$ be a parameter-independent program which, from each of its start histories, reaches one productive report. Its conditional expected attempt cost is at most $\kappa_+>0$ on every surviving history in every $\theta\in\Theta$. It uses $B$ live roles and returns to its start label on rejection; no reset of the physical state is assumed. It forms a virtual mark $x$, a current choice $J$ and a permitted export $T_Jr$ of its task readout. For every finite dictionary $C$ assume
\begin{equation}\label{eq:exports}
 |\{T_jr:j\in J,\ r\in C\}|\le D_0|C|,
 \qquad D_0\in\mathbb N.
\end{equation}
Every continuing phase and coin needed for this construction is included in the $B$ live roles. For fixed $0<l_+\le1\le u_+<\infty$ and $a_+,e_+\ge0$, assume for all Borel $q:X\to[0,1]$ and finite-image Borel $r:X\to H$, conditionally on every such history,
\begin{align}
 \E[q(x)]&\ge l_+\int q\,d\nu-a_+,\label{eq:upper-mass}\\
 \E[q(x)\|Z-T_Jr(x)\|^2]
 &\le u_+\int q(x)d_{r(x)}(x)\nu(dx)+e_+.
 \label{eq:upper-task}
\end{align}
The expectation here includes the within-attempt choice $J$. Unlike the lower, the upper is not required to hold for each choice separately. The choices retained after this call are nevertheless charged through the actual export or live role. These hypotheses express one common causal construction, not modelwise optimizers.

\section{An arithmetic envelope for exact memory at every phase}\label{sec:r71-arithmetic}
The spectral criterion of Theorem~\ref{thm:r69-spectral} determines exact stochastic memory by membership in the classical stochastic eigenvalue regions. A worst-phase bound and a matching estimate at a badly approximable fixed phase do not describe all fixed phases. In particular, excellent rational approximations can produce both unusually small memory and long intervening ranges of larger memory. We determine a single arithmetic envelope that controls the entire memory function, and derive its exact lower and upper logarithmic growth exponents for every irrational phase. The arithmetic is classical continued-fraction theory; the conclusions here concern its conversion into the least exact stochastic observer memory. We retain the distinction from the noise-dependent classical process bounds in \cite{Fanizza24} and from the classical irrationality-exponent identity in \cite{HoneVarona19}.

Throughout this section, $0<\epsilon<1$, $\alpha=1-\epsilon$, and
\[
 b_n=\tfrac12[1+\gamma\alpha^n\cos(2\pi nx+\varphi)],
 \qquad 0<\gamma\le1,\quad \sin(2\pi x)\ne0.
\]
The phase and nonzero contrast describe one prepared target, not information revealed to its observer. Write $m_x(\epsilon)=m(1-\epsilon,2\pi x)$ for its least live-label count, with a prescribed initial label, stochastic transitions and a decoder in $[0,1]$. Theorem~\ref{thm:r69-spectral} makes this independent of $\gamma,\varphi$. The actual squared-loss embedding of Theorem~\ref{thm:r67-embedding} forces exactly this probability response at the least physical cost $1/p$. Thus the conclusions below are about exact optimal implementation, not approximation of the target sequence or a new definition of loss.

For a positive integer $q$, put
\[
 e_x(q)=\frac{\|qx\|_{\mathbb R/\mathbb Z}}q
       =\min_{a\in\mathbb Z}|x-a/q|,
 \qquad
 \mathfrak A_x(\epsilon)
   =\min_{q\ge1}\max\{3,q,e_x(q)/\epsilon\}.
\]
The minimum is attained: its summand is at least $q$, while $q=1$ gives a finite upper. No global Diophantine lower hypothesis is imposed.

\begin{theorem}[Universal arithmetic envelope]\label{thm:r71-envelope}
For every nonreal phase and every $0<\epsilon<1$,
\begin{equation}\label{eq:r71-envelope}
 \tfrac12\mathfrak A_x(\epsilon)
 \le m_x(\epsilon)
 \le C_0\mathfrak A_x(\epsilon),\qquad C_0=1+4\pi^2.
\end{equation}
The constants are independent of the phase, contrast, preparation phase and damping. Each upper bound comes from an explicit stochastic circulant realization with a prescribed initial label and bounded decoder. The lower holds over all stochastic matrices and decoders, not only those circulants. The estimate is a constant-factor law for the full arithmetic envelope; it is not an assertion of an exact leading constant or a new classification of the stochastic spectral regions.
\end{theorem}
\begin{proof}
We use the polygon construction and the power-sector inequality in Lemma~\ref{lem:r69-arithmetic-tools}, whose positive-matrix proofs are included in Section~\ref{sec:r69-spectral}. Let $M=m_x(\epsilon)\ge3$. Among the $2M+1$ fractional parts of $0,x,\ldots,2Mx$, two lie in the same one of $2M$ intervals of length $1/(2M)$. Their difference supplies $1\le k\le2M$ with $\|kx\|\le1/(2M)$. If this norm is zero, using $q=k$ gives $\mathfrak A_x(\epsilon)\le2M$ directly.

Otherwise put $\delta=2\pi\|kx\|$ and $h=\pi/M$. Then $0<\delta\le h\le\pi/3$. For
$F=\cos\delta+\sin\delta\tan h$, one has $F\le\sec h\le2$ and
\[
 F-1\ge\tfrac34h\delta-\tfrac12\delta^2\ge h\delta/4.
\]
Here $\sin\delta\ge3\delta/4$, $\tan h\ge h$ and $1-\cos\delta\le\delta^2/2$ on the stated interval. The eigenvalue of the stochastic power $Q^k$ gives $(1-\epsilon)^kF\le1$. Bernoulli's inequality therefore yields
\[
 k\epsilon\ge1-(1-\epsilon)^k\ge1-1/F\ge h\delta/8.
\]
Consequently $e_x(k)/\epsilon=\delta/(2\pi k\epsilon)\le4M/\pi^2$. Since also $k\le2M$, the envelope is at most $2M$. This proves the lower.

For the upper fix any denominator $q$, choose its closest integer numerator $a$, and choose the smallest positive multiple $l$ of $q$ satisfying
\[
 l\ge3,\qquad l\ge4\pi^2e_x(q)/\epsilon.
\]
Then
$l<q+\max\{3,4\pi^2e_x(q)/\epsilon\}\le C_0\max\{3,q,e_x(q)/\epsilon\}$.
The rational angle $2\pi a/q$ is an $l$th-root angle. Thus the distance $\delta_l$ from $2\pi x$ to the nearest such root is at most $2\pi e_x(q)$. The directional polygon bound gives
\[
 1-R_l(2\pi x)\le2(\pi/l)\delta_l
             \le4\pi^2e_x(q)/l\le\epsilon.
\]
Hence $\alpha e^{2\pi ix}$ lies in the regular $l$-gon. The positive combination of its two adjacent roots and zero gives the stochastic matrix described in Lemma~\ref{lem:r69-arithmetic-tools}; its normalized eigenvector gives the bounded decoder from Theorem~\ref{thm:r69-spectral}. The original initial label suffices. Minimize over $q$ to conclude. All constructions act on the observer's private memory and do not add a physical observation or a call.
\end{proof}

For irrational $x$, $\mathfrak A_x(\epsilon)\to\infty$ as $\epsilon\downarrow0$: a bounded envelope would restrict $q$ to finitely many integers whose errors have a positive minimum. For a rational nonreal phase of reduced denominator $q$, the exact memory is eventually $q$ as $\epsilon\downarrow0$. Indeed $\alpha C_q+(1-\alpha)J$ gives the upper $q$, while a bounded sequence of smaller stochastic realizations would converge, by compactness, to one with the forbidden unit eigenvalue of order $q$. This rational boundary is different from assigning an irrationality exponent to a rational number. Zero contrast and real phases remain governed by their separately stated degeneracies.

\subsection{The arithmetic input, with proof}
For irrational $x$, changing it by an integer does not change any quantity above. Write its fractional part as $[0;a_1,a_2,\ldots]$, with convergents $p_n/q_n$. We use the standard initial values $p_{-1}=1$, $p_0=0$, $q_{-1}=0$, $q_0=1$ and the recursion
$p_{n+1}=a_{n+1}p_n+p_{n-1}$, $q_{n+1}=a_{n+1}q_n+q_{n-1}$. The possible initial repeated denominator is irrelevant below; all asymptotic statements use the eventually strictly increasing sequence.

\begin{lemma}[Convergent errors and the irrationality exponent]\label{lem:r71-cf}
Set $E_n=q_nx-p_n$, with $x$ reduced modulo one. Then
\begin{equation}\label{eq:r71-cf-error}
 \frac1{q_{n+1}+q_n}<|E_n|<\frac1{q_{n+1}},
 \qquad
 |qx-a|\ge|E_n|\quad(0<q<q_{n+1},\ a\in\mathbb Z).
\end{equation}
For all sufficiently large $n$, $e_x(q_n)=|E_n|/q_n$. Define
\[
 \mu(x)=\sup\{t:\ e_x(q)<q^{-t}\text{ for infinitely many }q\}.
\]
Then
\begin{equation}\label{eq:r71-mu}
 \mu(x)=1+\limsup_{n\to\infty}\frac{\log q_{n+1}}{\log q_n}
       \in[2,\infty].
\end{equation}
Moreover $x$ is badly approximable, in the sense that $e_x(q)\ge c/q^2$ for some $c>0$ and all $q$, if and only if the ratios $q_{n+1}/q_n$ are bounded.
\end{lemma}
\begin{proof}
The recursions give $p_nq_{n-1}-p_{n-1}q_n=(-1)^{n-1}$ by induction. If $t_{n+1}=[a_{n+1};a_{n+2},\ldots]$, repeated substitution in the continued fraction gives
\[
 x=\frac{p_nt_{n+1}+p_{n-1}}{q_nt_{n+1}+q_{n-1}},
 \quad |E_n|=\frac1{q_nt_{n+1}+q_{n-1}},
 \quad E_{n-1}=-t_{n+1}E_n.
\]
Since $a_{n+1}<t_{n+1}<a_{n+1}+1$, the denominator lies strictly between $q_{n+1}$ and $q_{n+1}+q_n$, proving the first pair of bounds. These errors eventually have modulus below $1/2$, so the numerator is the nearest integer.

For the second assertion, suppose $|qx-a|<|E_n|$. The determinant identity makes the two consecutive convergent vectors an integral basis. Write
$(a,q)=A(p_n,q_n)+B(p_{n-1},q_{n-1})$ with integers $A,B$.
Then $|A-Bt_{n+1}|<1$. If $B=0$, this forces $A=0$ and $q=0$. If $B\ge1$, it forces $A\ge Ba_{n+1}$ and hence $q\ge Bq_{n+1}\ge q_{n+1}$. If $B=-b\le-1$, it forces $A\le-ba_{n+1}$ and hence $q\le-bq_{n+1}<0$. These exhaust the possibilities and prove the asserted best-error inequality. No best-approximation assertion for arbitrary fitted rational data has been used.

Put $\Lambda=\limsup\log q_{n+1}/\log q_n\ge1$. The convergent upper $e_x(q_n)<1/(q_nq_{n+1})$ shows $\mu\ge1+\Lambda$, including the infinite case. If $s>\Lambda$ is finite, eventually $q_{n+1}\le q_n^s$. For $q_n\le q<q_{n+1}$, the second-kind inequality gives
\[
 e_x(q)>\frac1{q(q_{n+1}+q_n)}
        \ge\frac1{2q^{s+1}}.
\]
For every positive $\eta$, this excludes infinitely many errors below $q^{-(s+1+\eta)}$. Let $s\downarrow\Lambda$ and $\eta\downarrow0$ to prove \eqref{eq:r71-mu}.

If $q_{n+1}/q_n\le B$ eventually, the same interval argument gives $e_x(q)>1/[(B+1)q^2]$ there; reducing the positive constant handles the finitely many earlier denominators. Conversely $e_x(q_n)<1/(q_nq_{n+1})$ is incompatible with one positive lower constant when the ratios are unbounded. This proves the last equivalence. The identity \eqref{eq:r71-mu} is classical; it is also stated in the cited arithmetic literature \cite{HoneVarona19}, and is proved here to make the memory argument independent of that reference.
\end{proof}

\subsection{The complete logarithmic growth spectrum}
The envelope law retains more information than a single exponent. Its two logarithmic exponents nevertheless have a closed form for every irrational phase, including phases not known to satisfy a uniform approximation bound.

\begin{theorem}[Exact lower and upper memory exponents]\label{thm:r71-exponents}
For every irrational $x$, the exact stochastic memory obeys
\begin{align}
 \liminf_{\epsilon\downarrow0}
  \frac{\log m_x(\epsilon)}{\log(1/\epsilon)}
     &=\frac1{\mu(x)+1},\label{eq:r71-liminf}\\
 \limsup_{\epsilon\downarrow0}
  \frac{\log m_x(\epsilon)}{\log(1/\epsilon)}
     &=\frac{\mu(x)-1}{2\mu(x)-1}.\label{eq:r71-limsup}
\end{align}
For $\mu(x)=\infty$, the right sides mean $0$ and $1/2$. In particular the logarithmic exponent exists exactly when $\mu(x)=2$, and then it equals $1/3$. This does not imply a constant-factor cubic-root law for every such phase. The assertions describe one fixed phase while the noise tends to zero, not a worst-phase choice depending on the noise.
\end{theorem}
\begin{proof}
By Theorem~\ref{thm:r71-envelope}, multiplicative constants do not affect either logarithmic limit, so work with $\mathfrak A_x$.

First suppose $\mu<\infty$ and choose $t>\mu$. For all sufficiently large $q$, $e_x(q)\ge q^{-t}$. Therefore
\[
 \max\{q,e_x(q)/\epsilon\}
   \ge\max\{q,q^{-t}/\epsilon\}
   \ge\epsilon^{-1/(t+1)}.
\]
Every one of the finitely many remaining denominators has positive error, and its error term grows as $\epsilon^{-1}$, so those denominators do not invalidate this lower for sufficiently small $\epsilon$. This gives a lower limit at least $1/(\mu+1)$. Conversely, for $1<t<\mu$, infinitely many $q$ satisfy $e_x(q)<q^{-t}$. At $\epsilon=q^{-(t+1)}$, the envelope is at most $q$ once $q\ge3$. Taking $t\uparrow\mu$ proves \eqref{eq:r71-liminf}; for $\mu=\infty$ the same argument with arbitrarily large $t$, together with nonnegativity, gives zero.

For the upper limit with finite $\mu$, again choose $t>\mu$. Pigeonhole approximation at denominator bound $Q$ gives $1\le k\le Q$ with $e_x(k)\le1/(kQ)$. Such $k$ tend to infinity along $Q\to\infty$, because the phase is irrational. The eventual lower $e_x(k)\ge k^{-t}$ therefore forces $k\ge Q^{1/(t-1)}$. Hence
\[
 \mathfrak A_x(\epsilon)
 \le\max\{3,Q,\epsilon^{-1}Q^{-t/(t-1)}\}.
\]
Choose $Q=\lceil\epsilon^{-(t-1)/(2t-1)}\rceil$ and then let $t\downarrow\mu$. This bounds the upper limit by the right side of \eqref{eq:r71-limsup}. If $\mu=\infty$, use only $k\ge1$ and $Q=\lceil\epsilon^{-1/2}\rceil$ in the same pigeonhole argument to obtain the universal upper limit $1/2$.

To prove the reverse upper-limit bound, take consecutive large convergent denominators $Q=q_n$, $Q'=q_{n+1}$ and set
\begin{equation}\label{eq:r71-crossover}
 \epsilon_n=\frac1{Q(Q')^2}.
\end{equation}
For every rational $a/q$ with $0<q<Q'/2$, either $a/q=p_n/Q$, in which case its error is greater than $1/[Q(Q'+Q)]\ge1/(2QQ')$, or the determinant of the two integer pairs is nonzero and
\[
 |x-a/q|\ge\frac1{qQ}-|x-p_n/Q|
          >\frac1{qQ}-\frac1{QQ'}>\frac1{QQ'}.
\]
Thus all these denominators have $e_x(q)/\epsilon_n\ge Q'/2$. All other denominators have $q\ge Q'/2$. It follows that $\mathfrak A_x(\epsilon_n)\ge Q'/2$. On the other hand, denominator $Q'$ has
$e_x(Q')<1/(Q'q_{n+2})$, so
$e_x(Q')/\epsilon_n<Q Q'/q_{n+2}\le Q\le Q'$.
Consequently
\begin{equation}\label{eq:r71-crossover-envelope}
 Q'/2\le\mathfrak A_x(\epsilon_n)\le Q'.
\end{equation}
The corresponding logarithmic ratio is, up to a term tending to zero,
\[
 \frac{\log Q'}{\log Q+2\log Q'}
   =\frac{\log Q'/\log Q}{1+2\log Q'/\log Q}.
\]
The function on the right is increasing, continuous and tends to $1/2$ at infinity. Lemma~\ref{lem:r71-cf} proves precisely the remaining lower bound in \eqref{eq:r71-limsup}. Finally the two displayed exponents are equal exactly when $\mu=2$: solving the equality gives $\mu(\mu-2)=0$, and $\mu\ge2$.
\end{proof}

\begin{theorem}[Rigidity of the uniform cubic-root law]\label{thm:r71-cubic-rigidity}
For an irrational fixed phase, the uniform estimate
\[
 m_x(\epsilon)=\Theta(\epsilon^{-1/3})\qquad(\epsilon\downarrow0)
\]
holds if and only if $x$ is badly approximable. If it is not, the stronger oscillation statements hold:
\begin{equation}\label{eq:r71-oscillation}
 \liminf_{\epsilon\downarrow0}\epsilon^{1/3}m_x(\epsilon)=0,
 \qquad
 \limsup_{\epsilon\downarrow0}\epsilon^{1/3}m_x(\epsilon)=+\infty.
\end{equation}
Thus even an existing logarithmic exponent $1/3$ need not give uniform cubic-root constants.
\end{theorem}
\begin{proof}
For badly approximable $x$, Theorem~\ref{thm:r69-cubic} already proves both uniform bounds, and that result is inherited. Otherwise Lemma~\ref{lem:r71-cf} supplies consecutive denominators $Q,Q'$ with $Q'/Q\to\infty$ along a subsequence. At $\epsilon'_n=e_x(Q)/Q$, the envelope is at most $Q$, and the upper in \eqref{eq:r71-envelope} gives
\[
 (\epsilon'_n)^{1/3}m_x(\epsilon'_n)
 \le C_0(Q^2e_x(Q))^{1/3}
 <C_0(Q/Q')^{1/3}\longrightarrow0.
\]
At the crossover noise \eqref{eq:r71-crossover}, the lower in \eqref{eq:r71-envelope} and \eqref{eq:r71-crossover-envelope} instead give
\[
 \epsilon_n^{1/3}m_x(\epsilon_n)
 \ge\tfrac14(Q'/Q)^{1/3}\longrightarrow+\infty.
\]
Both noise sequences tend to zero. Nonnegativity proves the first equality in \eqref{eq:r71-oscillation}, and the second rules out a uniform upper constant. No convergence of sampled continued-fraction coefficients has been substituted for a global arithmetic property.
\end{proof}

These conclusions distinguish three questions: the complete envelope at a given noise, the logarithmic exponents determined by $\mu$, and uniform power-law constants governed by bounded convergent ratios. A quadratic irrational phase has bounded partial quotients and retains the inherited cubic-root law. For an explicit phase with exponent two but without uniform cubic constants, take $x=[0;1,2,3,4,\ldots]$. Its denominator ratios are unbounded, while the denominator recursion implies at least Fibonacci growth and
$\log(q_{n+1}/q_n)/\log q_n\le\log(n+2)/\log q_n\to0$.
Hence $\mu=2$ but \eqref{eq:r71-oscillation} holds. For $\mu=\infty$, the two exponents are $0$ and $1/2$. The stronger universal theorem does not determine the numerical value of $\mu$ for the distinct algebraic unit phase $(3+4i)/5$; no cubic-root claim for that phase is introduced.

\subsection{Several visible frequencies and leakage}
Let $r\ge1$, let $x_j\in(0,1/2)$ be distinct, and let $w_j>0$ with $\sum_jw_j\le1$. For prescribed phases $\varphi_j$, consider
\[
 b_n^{\rm mix}=\tfrac12\left[1+(1-\epsilon)^n
              \sum_{j=1}^r w_j\cos(2\pi nx_j+\varphi_j)\right].
\]
Let $m_{\rm mix}$ be its least exact live memory. For $0<\lambda<1$, let $s_{\rm mix}$ be the least memory of the decaying response $\lambda^n b_n^{\rm mix}$.

\begin{corollary}[Multi-frequency arithmetic envelopes]\label{cor:r71-mixtures}
For every $0<\epsilon<1$,
\begin{align}
 \max\{2r+1,\tfrac12\max_j\mathfrak A_{x_j}(\epsilon)\}
 &\le m_{\rm mix}(\epsilon)
 \le2+C_0\sum_j\mathfrak A_{x_j}(\epsilon),\label{eq:r71-mixture}\\
 \max\{2r+2,1+\tfrac12\max_j\mathfrak A_{x_j}(\epsilon)\}
 &\le s_{\rm mix}(\epsilon)
 \le3+C_0\sum_j\mathfrak A_{x_j}(\epsilon).\label{eq:r71-leaky-mixture}
\end{align}
In particular each count is comparable, with constants depending only on $r$, to the largest of the component envelopes. If at least one phase is irrational, their upper logarithmic exponent is
\[
 \max_{j:\,x_j\notin\mathbb Q}
       \frac{\mu(x_j)-1}{2\mu(x_j)-1},
\]
with the same infinite-exponent convention. We do not replace the lower limit of a maximum by the maximum of lower limits; good approximation scales for different phases need not coincide. For a single leaking frequency, the exact identity $s=1+m_x$ transfers both exponents and the uniform-law rigidity unchanged, independently of $\lambda$.

These are actual boundary-optimal squared-task observer counts under the retained embedding. The protected target can be realized on dimension at most $2r+1$ before leakage and at most $2r+2$ after leakage; the flag embedding doubles these sufficient physical dimensions. No optimal physical-dimension claim or uniform controller for unknown target phases is made.
\end{corollary}
\begin{proof}
The distinct modes $1,(1-\epsilon)e^{\pm2\pi ix_j}$ have nonzero coefficients. Their rank and individual spectral obstructions, together with the counted block-mixture construction, are exactly those of Theorem~\ref{thm:r69-finite-mixture}. Apply \eqref{eq:r71-envelope} to each component to obtain \eqref{eq:r71-mixture}.

For leakage, the $2r+1$ decaying scalar modes gain the independent constant mode in the complementary output, giving rank $2r+2$. The relative-pole obstruction of Proposition~\ref{prop:r70-relative-spectrum} requires $1+m(1-\epsilon,2\pi x_j)$ labels for each $j$. Conversely append one live zero-output state to the entire unattenuated stochastic realization and scale its continuing matrix by $\lambda$. This realizes the decaying mixture from the same initial label. These arguments prove \eqref{eq:r71-leaky-mixture}.

For a finite list, the upper logarithmic limit of its maximum is the maximum of its upper logarithmic limits. Rational components have bounded exact size and do not change this maximum when an irrational component is present. Theorem~\ref{thm:r71-exponents} proves the formula. The single-frequency assertion follows from the exact leakage increment of Theorem~\ref{thm:r70-leakage}. Finally the block-diagonal protected realization in Theorem~\ref{thm:r69-finite-mixture}, with one additional leakage state when needed, supplies the stated dimensions; the actual squared-loss and stopping-cost assertions follow from the same declared flag embedding. A private randomized block choice is stored in the counted labels and never reveals the protected block.
\end{proof}

The conclusions do not supply a sharp leading constant, a general exact intermediate-memory risk curve, or a minimal controller for an unspecified noncomputable phase. For a supplied computable phase and positive noise, interval bounds on $e_x(q)$ and the explicit polygon construction give finite certified upper constructions; exact minimality still uses the separate spectral-membership problem. Ideal finite labels, descriptions of transition probabilities, exact sampling and working storage remain distinct resources. The complete arithmetic envelope closes the fixed-phase constant-factor question for this response family, not the general multi-command foundations programme.

\section{Spectral and arithmetic laws of exact response memory}\label{sec:r69-spectral}
The instrument realization theorem preserves every terminal effect, whereas a prescribed task can require substantially less memory. We now determine the exact memory problem for a whole family of single-frequency protected responses, not just one damping factor or the previously solved three-label case. The criterion reduces the problem to classical stochastic eigenvalue regions. A self-contained sector argument and elementary rational approximation then give an exact worst-phase law and matching growth at every fixed badly approximable phase. None of the spectral-region, positive-realization, or sector machinery is claimed as newly invented \cite{DmitrievDynkin46,JohnsonPaparella17,KirklandSmigoc22}. The contributions considered here are their exact response-level equivalence with a prescribed initial label, their consequences for the physical squared task, and the phase-sensitive matched resource law.

\subsection{A physical family and the exact spectral criterion}
Let $0<\alpha\le1$, $0<\gamma\le1$, and $\vartheta,\varphi\in\mathbb R$, with $\sin\vartheta\ne0$. Consider the binary response
\begin{equation}\label{eq:r69-response}
 b_n=\frac12\bigl[1+\gamma\alpha^n\cos(n\vartheta+\varphi)\bigr],
 \qquad n\ge0.
\end{equation}
The rotation, contrast and preparation phase specify one prepared instance; independence of the minimum size does not assert one fixed decoder for unknown phases or contrasts. It is physically prepared on a qubit: take the Bloch vector
$\gamma(\sin\varphi,0,\cos\varphi)$, the channel
\[
 \Phi(X)=\alpha U_\vartheta XU_\vartheta^*
               +(1-\alpha)\tr(X)I/2,\qquad
 U_\vartheta=\begin{pmatrix}\cos(\vartheta/2)&-\sin(\vartheta/2)\\
                  \sin(\vartheta/2)&\cos(\vartheta/2)\end{pmatrix},
\]
and terminal effect $E=(I+\sigma_z)/2$. Its Born probabilities are exactly \eqref{eq:r69-response}. Theorem~\ref{thm:r67-embedding} embeds this response into the same three-outcome squared-loss experiment, with protected dimension four, fixed decoder dictionary $\{-1,1\}$, and least expected cost $1/p$. It forces the probability of output $+1$ after $n$ failures to equal $b_n$. Thus the mathematical realization problem below is the exact optimal-observer problem of an actual bounded squared task, not a sequence supplied by a loss oracle.

Write $m(\alpha,\vartheta;\gamma,\varphi)$ for the least $M$ for which
\[
 b_n=e_1^{\mathsf T}Q^nd\quad(n\ge0),\qquad
 Q\ge0,\quad Q\mathbf1=\mathbf1,\quad d\in[0,1]^M.
\]
The initial label is prescribed and counted. Let
\[
 \mathcal K_M=\{z\in\mathbb C:z\text{ is an eigenvalue of some }
                         M\times M\text{ row-stochastic matrix}\}.
\]
These are the classical Karpelevi\v c regions; their full boundary classification is prior work \cite{JohnsonPaparella17}. We will not use that classification as an unproved step in any bound below.

\begin{theorem}[Exact spectral realization, with no extra initial label]\label{thm:r69-spectral}
For the nondegenerate family \eqref{eq:r69-response},
\begin{equation}\label{eq:r69-minimal}
 m(\alpha,\vartheta;\gamma,\varphi)
       =\min\{M:\alpha e^{i\vartheta}\in\mathcal K_M\}.
\end{equation}
The value is independent of every nonzero contrast $\gamma$ and of the preparation phase $\varphi$; denote it by $m(\alpha,\vartheta)$. The criterion is necessary and sufficient over all stochastic matrices, including reducible matrices and arbitrary randomized decoders. For $0<\alpha<1$ it is finite. For $\alpha=1$ and $\vartheta/(2\pi)=a/q$ in lowest terms it is exactly $q$; for irrational $\vartheta/(2\pi)$ it is infinite.

In the squared-task embedding, these are exact live-label counts at cost $1/p$. At each fixed smaller memory bound, the attained risk is strictly above the full-history value $2$. Event-stopped and state-stopped total-label upper counts remain respectively $\max\{M,2\}$ and $M+2$.
\end{theorem}
\begin{proof}
Put $z=\alpha e^{i\vartheta}$. The response has the three distinct exponential modes $1,z,\bar z$, with nonzero coefficients. Cayley--Hamilton for any realizing $Q$ makes its characteristic polynomial an annihilating shift polynomial of $b_n$. Substitution and Vandermonde independence force that polynomial to vanish at $z$. This proves necessity without assuming a minimal linear representation has positive coordinates.

Conversely let $Qv=zv$ with $v\ne0$. Normalize a largest-modulus coordinate to one by multiplying $v$ by one complex scalar, and permute that coordinate into position one. Then $v_1=1$ and $|v_j|\le1$. The column
\[
 d_j=\tfrac12[1+\gamma\operatorname{Re}(e^{i\varphi}v_j)]
\]
lies in $[0,1]^M$. Row normalization and $Q^nv=z^nv$ give $e_1^{\mathsf T}Q^nd=b_n$ for every $n$. The same $M$ labels suffice: there is no preliminary random initialization or extra anchor. Unreachable labels can be pruned. Finiteness for $\alpha<1$ follows from the regular-polygon construction proved below, or from the inherited polygon upper in Theorem~\ref{thm:qubit-positive-phase}.

At $|z|=1$, equality in the triangle inequality at a largest-modulus eigenvector coordinate forces each positive successor to have the same modulus and value $z$ times the current value. Following positive successors reaches a directed cycle of length $l\le M$, so $z^l=1$. An irrational phase is therefore impossible. If $z$ has order $q$, every such cycle has length divisible by $q$, giving $M\ge q$; a $q$-cycle realizes it. Finally the squared embedding forces the whole response. Compact attainment at fixed $M$ was proved in Proposition~\ref{prop:r67-density}; equality of its risk with $2$ would give the excluded exact response. The existing packing lemma gives the terminal counts.
\end{proof}

The nondegeneracy assumptions cannot be suppressed. Zero contrast gives one constant label. At $\alpha=0$, $b_0=(1+\gamma\cos\varphi)/2$ and $b_n=1/2$ for $n\ge1$; one label suffices if $\gamma\cos\varphi=0$, and otherwise two are necessary and sufficient. Real phases and constant cases are classified in the supplement. In particular a vanishing oscillation amplitude is not the limit of an exact-memory lower valid for nonzero amplitude.

For algebraic $x=\operatorname{Re}z$, $y=\operatorname{Im}z\ne0$, spectral membership has the compact polynomial certificate
\begin{equation}\label{eq:r69-bilinear}
 \begin{gathered}
 Q\ge0,\ Q\mathbf1=\mathbf1,\quad
 Qa=xa-yc,\quad Qc=ya+xc,\\
 a_1=1,\ c_1=0,\quad a_j^2+c_j^2\le1\quad(1\le j\le M).
 \end{gathered}
\end{equation}
It has $M^2+2M$ real variables and polynomial degree at most two. Normalization of an eigenvector proves equivalence in both directions. Real-algebraic elimination decides it at each fixed $M$, and the finite polygon upper terminates a minimal-size search for $|z|<1$. This improves the form of the certificate, not the general complexity of real-algebraic elimination. A chosen algebraic preparation and contrast then yield algebraic decoder entries. State labels, encodings of those entries, exact-sampling resources and runtime remain separate accounts.

\subsection{The sector bound, its equality case, and exact finite phases}
We include the classical stochastic sector argument to expose every load-bearing inequality. The equality structure is consistent with the Type-0 classification of Kirkland and \v Smigoc \cite{KirklandSmigoc22}; it is not a new classification of stochastic spectral arcs.

\begin{lemma}[Stochastic sector inequality and rigidity]\label{lem:r69-sector}
If $z=x+iy$, $y\ne0$, is an eigenvalue of an $M\times M$ stochastic matrix, then $M\ge3$ and
\begin{equation}\label{eq:r69-sector}
 |y|\le(1-x)\cot(\pi/M).
\end{equation}
If equality holds, then, up to simultaneous relabelling and reversing the orientation when needed, that matrix is
\begin{equation}\label{eq:r69-lazy-cycle}
 Q=(1-t)I+tC_M,\qquad 0<t\le1,
\end{equation}
where $C_M$ is the directed $M$-cycle. Its eigenvalue on the equality ray is $1-t+te^{2\pi i/M}$, or its conjugate.
\end{lemma}
\begin{proof}
Conjugate the eigenvector when necessary so $y>0$. Let $P$ be the convex hull of its complex coordinates. The matrix
$e^{s(Q-I)}=e^{-s}\sum_{n\ge0}s^nQ^n/n!$ is stochastic for $s\ge0$, so multiplication by $e^{s(z-1)}$ maps $P$ into itself. Since $z$ is nonreal, $1-x>0$, and the spiral of any nonzero coordinate visits every angle, a finite set of points of this spiral surrounds zero. Hence zero lies in the interior of $P$.

List its $l\le M$ vertices counterclockwise as $r_je^{i\theta_j}$, with $r_j>0$. Their consecutive angular gaps $\Delta_j$ belong to $(0,\pi)$ and sum to $2\pi$. Put $u=1-x>0$ and $v=y>0$. The initial velocity $(-u+iv)r_je^{i\theta_j}$ must point into the half-plane to the left of the outgoing edge. Taking its signed two-dimensional cross product with that edge gives
\[
 r_j/r_{j+1}\le \cos\Delta_j+(u/v)\sin\Delta_j.
\]
All right sides are positive. Their product is at least one. Write $\phi=\arctan(u/v)\in(0,\pi/2)$. The function
$\log(\cos\Delta+(u/v)\sin\Delta)
 =\log\sqrt{1+(u/v)^2}+\log\cos(\Delta-\phi)$
is strictly concave on the interval of positive factors in $(0,\pi)$. Jensen's inequality therefore gives
\[
 1\le\prod_j[\cos\Delta_j+(u/v)\sin\Delta_j]
 \le[\cos(2\pi/l)+(u/v)\sin(2\pi/l)]^l.
\]
As $l\ge3$, this implies $u/v\ge\tan(\pi/l)\ge\tan(\pi/M)$, proving \eqref{eq:r69-sector}.

At equality, $l=M$, every gap is $2\pi/M$, and all radius ratios equal one: the polygon is regular and every coordinate is one distinct vertex. The velocity at a vertex is along its outgoing edge. But
$(z-1)v_j=\sum_kQ_{jk}(v_k-v_j)$ is a nonnegative combination of chords. Strict convexity of this polygon makes all chords other than the self-chord and outgoing edge point strictly to the interior side of that edge. Their coefficients must therefore vanish. The remaining coefficient $t=Q_{j,j+1}$ satisfies $z-1=t(e^{2\pi i/M}-1)$ and is independent of $j$. Row normalization gives \eqref{eq:r69-lazy-cycle} and $t\le1$.
\end{proof}

\begin{theorem}[Exact first-sector memory threshold and optimizer]\label{thm:r69-sector-phase}
For $M\ge3$ and $0<\vartheta\le2\pi/M$, a response realization with at most $M$ live labels exists exactly when
\begin{equation}\label{eq:r69-sector-threshold}
 0<\alpha\le r_M(\vartheta)
 :=\frac1{\cos\vartheta+\sin\vartheta\tan(\pi/M)}.
\end{equation}
The threshold is attained. At $\alpha=r_M(\vartheta)$ the least memory is exactly $M$, the transition matrix is the lazy directed cycle in \eqref{eq:r69-lazy-cycle}, and, in cyclic coordinates with prescribed initial label zero, its decoder is uniquely
\[
 d_j=\tfrac12[1+\gamma\cos(\varphi+2\pi j/M)].
\]
This is global over arbitrary stochastic transition and decoder rows, not an optimization restricted in advance to circulants.
\end{theorem}
\begin{proof}
The spectral criterion and \eqref{eq:r69-sector} prove necessity. For sufficiency put $\omega=e^{2\pi i/M}$ and $J=\mathbf1\mathbf1^{\mathsf T}/M$. The nonnegative coefficients
\[
 b=\frac{\alpha\sin\vartheta}{\sin(2\pi/M)},\qquad
 a=\alpha\cos\vartheta-b\cos(2\pi/M),\qquad
 c=1-a-b
\]
sum to one under \eqref{eq:r69-sector-threshold}. The inequalities $a,b\ge0$ follow from $0<\vartheta\le2\pi/M$; $c\ge0$ is exactly the stated threshold. Then $Q=aI+bC_M+cJ$ has eigenvector $(1,\omega,\ldots,\omega^{M-1})$ at eigenvalue $\alpha e^{i\vartheta}$ and the displayed decoder realizes the response.

At the threshold, strict monotonicity of $\cot(\pi/M)$ excludes every smaller size. Lemma~\ref{lem:r69-sector} fixes $Q$. Here $b=t>0$, so $C_M=(Q-(1-t)I)/t$. The rows $e_0^{\mathsf T}Q^n$, $0\le n<M$, span all coordinate rows, because the rows $e_0^{\mathsf T}C_M^j$ do and each is a polynomial combination of the former. Thus all response coefficients determine the decoder uniquely. This proves the rigidity assertion including the endpoint $t=1$.
\end{proof}

For the inherited algebraic angle $\cos\vartheta=3/5$, $\sin\vartheta=4/5$, let
\[
 a_j=\frac5{3+4\tan(\pi/j)},\qquad j=3,4,5,6.
\]
One has $0<\vartheta<2\pi/6$. Consequently exact memory is three on $(0,a_3]$, four on $(a_3,a_4]$, five on $(a_4,a_5]$, and six on $(a_5,a_6]$. In particular $a_4=5/7$. Above $a_6$ at least seven labels are needed; no assertion that a regular polygon is always globally optimal beyond the first sector is made. The old three-label threshold is recovered rather than reproved as a new independent result.

\subsection{Exact worst-phase size and arithmetic improvements}
For an integer $q\ge3$ define the distance to the nearest $q$th root angle by
\[
 \delta_q(\vartheta)=\min_{j\in\mathbb Z}|\vartheta-2\pi j/q|
             \quad\text{modulo }2\pi,\qquad 0\le\delta_q\le\pi/q,
\]
and define the radial boundary of the regular polygon by
\begin{equation}\label{eq:r69-polygon-radius}
 R_q(\vartheta)=\frac{\cos(\pi/q)}{\cos(\pi/q-\delta_q(\vartheta))}.
\end{equation}

\begin{lemma}[Directional polygon construction and power obstruction]\label{lem:r69-arithmetic-tools}
If $0<\alpha\le R_q(\vartheta)$, a $q$-label exact response realization exists. Conversely, any $M$-label realization satisfies, for every positive integer $k$,
\begin{equation}\label{eq:r69-power-obstruction}
 \alpha^k\bigl[\cos(k\vartheta)+|\sin(k\vartheta)|\tan(\pi/M)\bigr]\le1.
\end{equation}
For $h=\pi/q$ and $0\le\delta=\delta_q\le h$,
\begin{equation}\label{eq:r69-polygon-error}
 1-R_q(\vartheta)\le2h\delta.
\end{equation}
\end{lemma}
\begin{proof}
The intersection of the ray at angle $\vartheta$ with the chord joining its two adjacent unit roots has radius \eqref{eq:r69-polygon-radius}. Thus $\alpha e^{i\vartheta}$ is in their triangle with zero. Write it as $a\omega^j+b\omega^{j+1}+c0$, with $a,b,c\ge0$ and sum one. The stochastic matrix $aC_q^j+bC_q^{j+1}+cJ$ has that eigenvalue, giving the first assertion by Theorem~\ref{thm:r69-spectral}. The matrix $Q^k$ of any realization is still stochastic and has eigenvalue $\alpha^ke^{ik\vartheta}$. Apply the sector lemma; for a real eigenvalue the displayed inequality is automatic. Finally $\cos(h-\delta)\ge\cos h\ge1/2$ and
$\cos(h-\delta)-\cos h=2\sin(h-\delta/2)\sin(\delta/2)\le h\delta$.
Division proves \eqref{eq:r69-polygon-error}.
\end{proof}

\begin{theorem}[Exact worst-phase memory]\label{thm:r69-worst-phase}
For every $0<\alpha<1$,
\begin{equation}\label{eq:r69-worst-phase}
 \max_{\sin\vartheta\ne0}m(\alpha,\vartheta)
  =\max\left\{3,\left\lceil\frac\pi{\arccos\alpha}\right\rceil\right\}.
\end{equation}
In particular the worst-phase exact label count is asymptotic to
$\pi/\sqrt{2(1-\alpha)}$ as $\alpha\uparrow1$. The phase may depend on $\alpha$; this is not a lower rate for each fixed phase.
\end{theorem}
\begin{proof}
Let the displayed integer be $K$. The regular $K$-gon contains the disk of radius $\cos(\pi/K)\ge\alpha$, so the polygon lemma gives the upper simultaneously for all phases. If $K=3$, three distinct nonzero modes force at least three states. Otherwise put $M=K-1\ge3$ and choose $\vartheta=\pi/M$. The definition of $K$ gives $\alpha>\cos(\pi/M)$; the exact first-sector threshold at this angle is $\cos(\pi/M)$. No $M$-label realization exists, proving the lower and attainment of the maximum. The asymptotic follows from $1-\cos t\sim t^2/2$.
\end{proof}

A number $x$ is called badly approximable here with constant $c>0$ when
\begin{equation}\label{eq:r69-diophantine}
 \|kx\|_{\mathbb R/\mathbb Z}\ge c/k\qquad(k\ge1).
\end{equation}
The constant is part of the hypothesis, not something inferred from a finite list of continued-fraction coefficients. The next theorem holds for every such phase and proves a different, matching fixed-phase rate.

\begin{theorem}[Sharp cubic-root memory at a fixed Diophantine phase]\label{thm:r69-cubic}
Let $x=\vartheta/(2\pi)$ satisfy \eqref{eq:r69-diophantine}. Then, for $0<\alpha<1$ and $\epsilon=1-\alpha$,
\begin{equation}\label{eq:r69-cubic-bounds}
 \left(\frac{\pi^2c}{16\epsilon}\right)^{1/3}
 \le m(\alpha,\vartheta)
 \le \max\left\{4,\left\lceil\frac3c\right\rceil,
              \left\lceil\left(\frac{4\pi^2}{c^2\epsilon}\right)^{1/3}\right\rceil\right\}.
\end{equation}
Thus $m(\alpha,\vartheta)=\Theta_c((1-\alpha)^{-1/3})$. The exact live-state bits are $\frac13\log_2(1/(1-\alpha))+O_c(1)$ in the ideal finite-label model. These are exact stochastic realizations at the original least expected-cost budget, not approximation of the response or additional physical observations. The statement does not supply matching program-bit or exact-sampling complexity bounds.
\end{theorem}
\begin{proof}
We use only the following elementary approximation fact. For any irrational $x$ and integer $Q\ge1$, among the $Q+1$ fractional parts $0,x,\ldots,Qx$ two fall in the same one of $Q$ equal intervals. Subtracting them gives $1\le k\le Q$ with $0<\|kx\|\le1/Q$.

For the lower, take a realizing size $M\ge3$ and apply that fact with $Q=2M$. Set $\delta=2\pi\|kx\|$ and $h=\pi/M$. Then
\[
 \pi c/M\le\delta\le h\le\pi/3.
\]
Writing $F=\cos\delta+\sin\delta\tan h$, we have $F\le\sec h\le2$. Also $\sin\delta\ge3\delta/4$ on $[0,\pi/3]$, $\tan h\ge h$ and $1-\cos\delta\le\delta^2/2$, so
\[
 F-1\ge3h\delta/4-\delta^2/2\ge h\delta/4.
\]
By \eqref{eq:r69-power-obstruction}, $\alpha^kF\le1$. Hence
\[
 2M\epsilon\ge k\epsilon\ge1-\alpha^k
       \ge1-1/F\ge h\delta/8
       \ge\pi^2c/(8M^2).
\]
This proves the lower constant.

For the upper, let $Q$ be the integer on the right of \eqref{eq:r69-cubic-bounds}. The approximation fact provides integers $1\le q\le Q$ and $a$ with $|qx-a|\le1/Q$. The lower bound \eqref{eq:r69-diophantine} forces $q\ge cQ\ge3$. The distance from $\vartheta$ to a $q$th-root angle is at most $2\pi/(qQ)$. The polygon-error bound gives
\[
 1-R_q(\vartheta)\le\frac{4\pi^2}{q^2Q}
       \le\frac{4\pi^2}{c^2Q^3}\le\epsilon.
\]
Thus $\alpha\le R_q(\vartheta)$ and the explicit circulant construction realizes the response with $q\le Q$ labels. Both bounds apply for the same fixed phase as $\epsilon\downarrow0$. Taking base-two logarithms and rounding the label bits proves the last assertion.
\end{proof}

The hypothesis is nonempty and has elementary physical examples. For
$x=(3-\sqrt5)/2$, choose the closest integer $a$ to $kx$. The integer
$a^2-3ak+k^2$ is nonzero. Since the other root is $x'=(3+\sqrt5)/2$ and
$|a/k-x'|\le\sqrt5+1/2<3$, factorization gives
$|a/k-x|>1/(3k^2)$. Thus \eqref{eq:r69-diophantine} holds with $c=1/3$. This provides a fixed qubit rotation with the matched exponent. Its trigonometric primitive entries are not asserted to be algebraic. In particular the Diophantine condition has not been proved here for the different algebraic rotation $(3+4i)/5$.

\begin{corollary}[A stronger fixed-algebraic-angle exclusion]\label{cor:r69-algebraic-angle}
For the inherited angle $e^{i\vartheta}=(3+4i)/5$ and $0<\alpha<1$, every exact $M$-label stochastic realization satisfies
\begin{equation}\label{eq:r69-algebraic-bound}
 1-\alpha\ge\frac3{16M^2\,25^M}.
\end{equation}
Consequently its exact memory is at least of order $\log(1/(1-\alpha))$ as $\alpha\uparrow1$. This improves the earlier $\log/\log\log$ lower, without assigning the cubic-root law to this angle or claiming a matching upper for it.
\end{corollary}
\begin{proof}
Again choose $1\le k\le2M$ with $\delta=2\pi\|k\vartheta/(2\pi)\|\le\pi/M$. The Gaussian integer $(3+4i)^k-5^k$ is nonzero: otherwise the unit complex number would be a root of unity, forcing $6/5$ to be a rational algebraic integer. Its modulus is at least one. Therefore
$\delta\ge|1-e^{ik\vartheta}|\ge5^{-k}\ge25^{-M}$.
Repeat the lower proof of Theorem~\ref{thm:r69-cubic} up to $2M(1-\alpha)\ge h\delta/8$, with $h=\pi/M$. This gives the stronger constant $\pi/(16M^2 25^M)$ and hence \eqref{eq:r69-algebraic-bound}. Taking logarithms gives
$\log(1/(1-\alpha))\le\log(16/3)+2\log M+M\log25$, proving the order assertion.
\end{proof}

\subsection{Finite spectral mixtures and higher-dimensional protected systems}
The matching exponent is not confined to a two-dimensional physical carrier. Let $r\ge1$, let $\vartheta_1,\ldots,\vartheta_r\in(0,\pi)$ be distinct, and let $w_j>0$ with $\sum_jw_j\le1$. For fixed phases $\varphi_j$, consider
\begin{equation}\label{eq:r69-mixture}
 b_n^{\rm mix}=\tfrac12\left[1+\alpha^n
             \sum_{j=1}^r w_j\cos(n\vartheta_j+\varphi_j)\right].
\end{equation}
Write $m_{\rm mix}(\alpha)$ for its least pure-start stochastic memory.

\begin{theorem}[Spectral mixtures preserve the matched arithmetic exponent]\label{thm:r69-finite-mixture}
For $0<\alpha<1$, the exact memory of \eqref{eq:r69-mixture} satisfies
\begin{equation}\label{eq:r69-mixture-bounds}
 \max\{2r+1,\max_j m(\alpha,\vartheta_j)\}
 \le m_{\rm mix}(\alpha)
 \le 1+\one_{\{\sum_jw_j<1\}}+\sum_{j=1}^r m(\alpha,\vartheta_j).
\end{equation}
All models are realized with the counted prescribed initial label. If every fixed $\vartheta_j/(2\pi)$ has a global badly-approximable certificate, then
$m_{\rm mix}(\alpha)=\Theta((1-\alpha)^{-1/3})$, with constants depending on the finite phase list. The same conclusion holds when some phases are rational and at least one is certified badly approximable. This probability response is physically prepared on dimension at most $2r+1$, and its exact memory is the boundary-optimal observer memory of the same squared task after the standard two-dimensional herald-flag embedding.
\end{theorem}
\begin{proof}
The $2r+1$ modes $1,\alpha e^{\pm i\vartheta_j}$ are pairwise distinct and have nonzero coefficients. Their Hankel rank is $2r+1$ by Vandermonde independence. Every realizing characteristic polynomial vanishes at every nonconstant mode, giving both lower bounds through Theorem~\ref{thm:r69-spectral}.

For each component choose a minimal stochastic realization $(Q_j,d_j)$ starting in its first label. Keep their state sets disjoint and add one initial label $o$ with decoder $b_0^{\rm mix}$. On its first failure, its row enters block $j$ with subprobability row $w_je_1^{\mathsf T}Q_j$. If $w_0=1-\sum_jw_j>0$, its remaining mass enters one absorbing label with decoder $1/2$. Within each block use $Q_j$. Thus after $n\ge1$ failures the output probability is
$w_0/2+\sum_jw_je_1^{\mathsf T}Q_j^nd_j=b_n^{\rm mix}$, while the initial decoder handles $n=0$. No physical call is added and no persistent choice is stored outside the counted labels. This proves the upper. Theorem~\ref{thm:r69-cubic} supplies matched exponents for the finite certified phase list. For a rational phase of order $q$, the matrix $\alpha C_q+(1-\alpha)J$ gives a uniform $q$-label upper, so adding those components does not change the exponent.

For the physical realization, take a direct sum of $r$ qubits, prepare the phase-rotated pure state in block $j$ with weight $w_j$, and use the damped rotation channel in that block. Drop off-diagonal block coherences. If $w_0>0$, add one scalar block of mass $w_0$ with effect $1/2$ and identity evolution. Each diagonal branch is completely positive and their direct sum is trace preserving. The terminal effect is the direct sum of the qubit $Z$ effects and the scalar half effect. Its probability is exactly \eqref{eq:r69-mixture}. Theorem~\ref{thm:r67-embedding} therefore gives the same fixed squared loss and cost $1/p$, and forces this response at the optimum. The observer's independently randomized block selection does not reveal or copy the protected block.
\end{proof}

The three memory statements have different quantifiers: the exact worst-phase law permits the phase to vary with the noise, the cubic-root law fixes any phase with the displayed arithmetic certificate, and the last corollary fixes the older algebraic rotation without such a certificate. Noise-dependent classical memory bounds for quantum processes were already established by Fanizza, Lumbreras and Winter \cite{Fanizza24}; their worst-phase process bound is not a fixed-phase matched law for this response task. The stochastic regions and their Type-0 rigidity are also classical. The present proofs make the task-level reduction and the new arithmetic comparison explicit, rather than claiming earliest priority from a literature search. No statement in this section solves the different squared-sign prediction frontier or the phase-type timing frontier.

\section{Transient output faces and leakage-sensitive memory}\label{sec:r70-leakage}
Uniform decay of a physical state and decay of its observable probabilities do not measure the same obstruction to exact stochastic realization. In this section leakage removes population from a protected subspace, whereas internal depolarization changes the spectrum relative to that surviving population. We prove that the first operation costs exactly one additional live label in the nondegenerate single-frequency family, independently of how rapid the leakage is. We then obtain finite-model witnesses on a fixed six-dimensional squared-prediction experiment, and an exact finite-realization result when a positive background is added to the terminal effect.

Nonnegative realizations, dominant-pole obstructions and transient-state reductions have classical antecedents \cite{BenvenutiFarina04,OCinneide90,Maier91,MonrasWinter14}. We do not claim to originate these methods. The statements below specify the constrained stochastic decoder, prescribed initial label, sharp label increment, physical leakage channel and finite squared-task tests. All necessary positive-matrix arguments are included.

\subsection{The output-visible transient core}
For a nonzero sequence $u=(u_n)_{n\ge0}$ in $[0,1]$ tending to zero, let $s(u)$ be its least stochastic realization size:
\[
 u_n=e_1^{\mathsf T}Q^nd,\qquad Q\ge0,\quad Q\one=\one,\quad 0\le d\le\one.
\]
The initial label is fixed, not a freely randomized initial distribution. Define $t(u)$ by the same formula with a nonnegative row-substochastic matrix $S$, $S\one\le\one$, whose spectral radius is strictly below one. Both minima are infinite if no finite realization exists.

\begin{lemma}[Exact transient-core reduction]\label{lem:r70-core}
For every such sequence,
\begin{equation}\label{eq:r70-core-count}
 s(u)=t(u)+1.
\end{equation}
More precisely, any $M$-label stochastic realization has a nonempty set of at most $M-1$ reachable labels that can eventually give positive output. Its restriction $S$ realizes $u$, is transient and uses the same initial label. If also $u_n\le C\lambda^n$ for $0<\lambda<1$, then the spectral radius of this restriction is at most $\lambda$.
\end{lemma}
\begin{proof}
Remove labels not reachable from label one. Call a remaining label active if a positive-probability path from it reaches some label $j$ with $d_j>0$. The initial label is active because $u$ is nonzero. A path that leaves the active set cannot return: a return would make every intervening reachable label active. Thus, with $S$ the active principal submatrix and $d_A$ its decoder, $u_n=e_1^{\mathsf T}S^nd_A$.

For each active $i$, choose a path of length $a_i$ from the start to $i$ of weight $p_i>0$. For each active $j$, choose a path of length $b_j$ from $j$ with positive terminal expected output $r_j>0$. Positivity gives, for every $n$,
\begin{equation}\label{eq:r70-path-bound}
 p_ir_j(S^n)_{ij}\le u_{a_i+n+b_j}.
\end{equation}
Consequently every entry of $S^n$ tends to zero. In finite dimension this forces every eigenvalue to have modulus less than one. In the exponential case it gives $(S^n)_{ij}\le C_{ij}\lambda^n$, and applying this bound to an eigenvector proves the stronger spectral-radius bound. There must be a reachable inactive label: otherwise $S$ would be stochastic and have eigenvalue one. Hence the active size is at most $M-1$, proving $t(u)+1\le s(u)$.

Conversely a $k$-label transient realization extends to the stochastic matrix and decoder
\begin{equation}\label{eq:r70-cemetery}
 \widehat Q=\begin{pmatrix}S&\one-S\one\\0&1\end{pmatrix},
 \qquad \widehat d=\binom{d_A}{0}.
\end{equation}
They realize the same sequence from the same initial label with $k+1$ labels. This proves the reverse inequality, including the infinite case. The added label is a live zero-output label until the actual herald; it is not permission to stop at an unobserved physical leakage event.
\end{proof}

We recall a small positive-matrix fact with its proof. If $B\ge0$ is irreducible and has a nonzero eigenvalue, it has a vector $h>0$ with $Bh=Rh$, where $R$ is its spectral radius. Indeed, for $B+\epsilon\one\one^{\mathsf T}$ the map sending a simplex vector to its normalized image has a fixed point by the finite-dimensional simplex fixed-point theorem. The fixed point is strictly positive. Along $\epsilon\downarrow0$ a convergent subsequence gives $Bh=Rh$, $h\ge0$, $\sum_i h_i=1$. If a coordinate vanished, the eigenvector equation and irreducibility would propagate that zero to every coordinate. Thus $h>0$, and $R>0$. The similarity $R^{-1}\operatorname{diag}(h)^{-1}B\operatorname{diag}(h)$ is stochastic, so its maximum-row-sum bound shows that $R$ is indeed the spectral radius. This is the Perron normalization used below, not a new Perron theorem.

\begin{proposition}[A relative spectral obstruction for decaying outputs]\label{prop:r70-relative-spectrum}
Suppose an $M$-label stochastic response satisfies $0\le u_n\le C\lambda^n$ with $0<\lambda<1$. If its rational generating function has a nonreal pole at $\zeta^{-1}$, then
\begin{equation}\label{eq:r70-relative-spectrum}
 \zeta/\lambda\in\mathcal K_{M-1}.
\end{equation}
Here $\mathcal K_j$ is the stochastic eigenvalue region defined in Section~\ref{sec:r69-spectral}. It is the decay-normalized mode, not merely the absolute mode $\zeta$, that constrains the active state count.
\end{proposition}
\begin{proof}
The core in Lemma~\ref{lem:r70-core} gives
$\sum_nu_nz^n=e_1^{\mathsf T}(I-zS)^{-1}d_A$. Hence $\zeta$ is an eigenvalue of $S$. Order its strongly connected components to obtain a block triangular matrix. One irreducible diagonal block $B$ has that eigenvalue. Its spectral radius $R$ obeys $|\zeta|\le R\le\lambda$. Perron normalization gives a stochastic matrix $T$ of size $l\le M-1$ with eigenvalue $\zeta/R$ and eigenvector $v$.

A stationary probability row $\pi$ for $T$ exists: take a convergent subsequence of the averages of any probability row and its first $n-1$ iterates. Since $\zeta/R\ne1$, one has $\pi v=0$. Put $a=R/\lambda\in(0,1]$. The stochastic matrix
$T'=aT+(1-a)\one\pi$ satisfies $T'v=(\zeta/\lambda)v$. Padding by unreachable self-loop labels, if needed, proves membership in $\mathcal K_{M-1}$. No global Perron scaling of the reducible core has been assumed.
\end{proof}

The last qualification matters. A reducible matrix such as
$S=\left(\begin{smallmatrix}1/2&1/4\\0&1/4\end{smallmatrix}\right)$
has spectral radius $1/2$, but no strictly positive $h$ satisfies $Sh\le h/2$. The proof normalizes the irreducible block carrying the observed pole, rather than silently imposing that false global inequality.

\subsection{An exact leakage law for every damping and phase}
Fix $0<\lambda<1$, $0<\alpha\le1$, $0<\gamma\le1$, and $\sin\vartheta\ne0$. Let
\begin{equation}\label{eq:r70-leaky-response}
 u_n=\frac{\lambda^n}{2}
       [1+\gamma\alpha^n\cos(n\vartheta+\varphi)].
\end{equation}
The scalar $\lambda$ is survival against leakage; $\alpha$ is coherence retention within the surviving qubit. These parameters are known properties of the prepared target channel, not extra observations available to the observer.

\begin{theorem}[Exact leakage increment]\label{thm:r70-leakage}
For the nondegenerate family \eqref{eq:r70-leaky-response},
\begin{equation}\label{eq:r70-exact}
 s(u)=1+m(\alpha,\vartheta)
     =1+\min\{k:\alpha e^{i\vartheta}\in\mathcal K_k\}.
\end{equation}
This holds over all stochastic matrices and randomized decoder rows, with a prescribed initial label. In particular the value is independent of every $0<\lambda<1$, every nonzero contrast and every preparation phase. For $\alpha<1$ it is finite. At $\alpha=1$ it is $h+1$ if $e^{i\vartheta}$ has finite order $h$, and it is infinite otherwise.

The full two-output Hankel response has rank four, regardless of this possibly much larger or infinite exact memory. In the fixed squared-task embedding it is the least live memory at the original cost $1/p$, with event-stopped and state-stopped total-label upper counts $\max\{s(u),2\}$ and $s(u)+2$ when finite.
\end{theorem}
\begin{proof}
The three modes $\lambda$, $\lambda\alpha e^{i\vartheta}$ and $\lambda\alpha e^{-i\vartheta}$ in $u$ are distinct and have nonzero coefficients. Thus its generating function has a pole at $(\lambda\alpha e^{i\vartheta})^{-1}$. Proposition~\ref{prop:r70-relative-spectrum} makes every finite realizing size $M$ satisfy $\alpha e^{i\vartheta}\in\mathcal K_{M-1}$. Theorem~\ref{thm:r69-spectral} proves the required lower and the infinite case.

For the upper take a minimal realization $(Q,d)$ of the unattenuated response
$b_n=[1+\gamma\alpha^n\cos(n\vartheta+\varphi)]/2$.
Replace its continuing block by $\lambda Q$ and append a live zero-output label, as in \eqref{eq:r70-cemetery}. Its output is exactly $\lambda^nb_n$ from time zero. This proves equality without an extra initialization operation. The unit-circle alternatives follow from the already proved root-of-unity characterization.

The scalar response spans its three exponential modes. Its complement adds the independent constant mode one, because $\lambda<1$. Vandermonde independence therefore gives full response rank four. Theorem~\ref{thm:r67-embedding} forces this entire probability response in the actual squared task; the role-packing statement is then inherited without identifying protected physical states with live observer labels.
\end{proof}

\begin{corollary}[Noise laws depend on relative damping]\label{cor:r70-rates}
For fixed $0<\lambda<1$ and $0<\alpha<1$, the worst-phase exact count is
\[
 \max_{\sin\vartheta\ne0}s(u)
 =1+\max\{3,\lceil\pi/\arccos\alpha\rceil\}.
\]
For every fixed phase with the global badly-approximable certificate of Theorem~\ref{thm:r69-cubic}, $s(u)-1$ has the same matching order $(1-\alpha)^{-1/3}$. Neither rate depends on leakage $\lambda$, even when leakage is allowed to vary simultaneously with $\alpha$. At the original algebraic phase, the first nontrivial threshold is exactly four live labels for $0<\alpha\le5/(3+4\sqrt3)$, and at least five beyond it. No badly-approximable claim is made for that algebraic phase.
\end{corollary}
\begin{proof}
Substitute \eqref{eq:r70-exact} into the exact worst-phase, fixed-phase and three-label results. All subtractions concern live labels, including the necessary zero-output label; no terminal label has been subtracted from an event-stopped dictionary.
\end{proof}

\subsection{A strictly contractive quantum channel with infinite exact memory}
Let $K=\mathbb C^2\oplus\mathbb C|\bot\rangle$, let $P$ project onto the qubit and let $\Omega=|\bot\rangle\langle\bot|$. With the qubit channel $\Phi_\alpha$ and preparation from Section~\ref{sec:r69-spectral}, define
\begin{equation}\label{eq:r70-channel}
 \Psi_{\lambda,\alpha}(X)
 =\lambda\Phi_\alpha(PXP)
  +[\tr((I-P)X)+(1-\lambda)\tr(PX)]\Omega.
\end{equation}
The qubit output is embedded in its original block and cross-block coherences are discarded. Prepare the qubit and audit the qubit effect $E=(I+\sigma_z)/2$, extended by zero at $\bot$.

\begin{theorem}[Uniform physical mixing versus exact response memory]\label{thm:r70-physical}
Equation \eqref{eq:r70-channel} is a completely positive trace-preserving channel with unique stationary state $\Omega$. For all density matrices $\rho,\sigma$ and all $n\ge0$,
\[
 \|\Psi_{\lambda,\alpha}(\rho)-\Psi_{\lambda,\alpha}(\sigma)\|_1
 \le\lambda\|\rho-\sigma\|_1,
 \qquad
 \|\Psi_{\lambda,\alpha}^n(\rho)-\Omega\|_1\le2\lambda^n.
\]
Its prepared audit response is exactly \eqref{eq:r70-leaky-response}, so its exact stochastic observer memory is \eqref{eq:r70-exact}. In particular, at $\alpha=1$ and the algebraic rotation $(3+4i)/5$, every $0<\lambda<1$ gives infinite exact memory despite strict trace-norm contraction and exponential convergence to one stationary state. For algebraic $\lambda$ and algebraic preparation data, all primitive superoperator entries are algebraic.

The squared-prediction embedding has protected dimension six, fixed dictionary $\{-1,1\}$, three actual audit outcomes and least expected cost $1/p$. In the specified irrational, internally undamped case, every finite live-memory class has an attained strictly larger risk than the full-history value $2$. A prefix with $T$ exact-output labels and one live zero-output tail label has the same cost and excess risk at most
\begin{equation}\label{eq:r70-approx}
 \frac{2\varepsilon p(q\bar t\lambda)^T}{1-q\bar t\lambda},
 \qquad q=1-p,\quad\varepsilon=1/4.
\end{equation}
\end{theorem}
\begin{proof}
Write $\Xi(X)=\Phi_\alpha(PXP)+\tr((I-P)X)\Omega$ and
$\mathcal E(X)=\tr(X)\Omega$. Both are channels: compression, the declared qubit channel and scalar replacement are completely positive, and their output traces sum to the input trace. Then $\Psi_{\lambda,\alpha}=\lambda\Xi+(1-\lambda)\mathcal E$. On a traceless Hermitian difference the second term vanishes. Trace-norm contraction of a channel follows directly by applying positivity and trace preservation to the positive and negative parts of that difference. This proves the contraction inequality. Since $\Omega$ is fixed, iteration gives the displayed bound and uniqueness.

Starting in the qubit, its surviving block after $n$ calls is $\lambda^n\Phi_\alpha^n(\rho)$; the complementary mass is at $\bot$. Taking the effect gives \eqref{eq:r70-leaky-response}. The preceding theorem proves the memory claim. At $(3+4i)/5$ the unit phase is not a root of unity, by the rational algebraic-integer argument already proved. The spectral obstruction is in the decaying output-visible block, not at the peripheral spectrum of the entire trace-preserving channel.

The embedding of Theorem~\ref{thm:r67-embedding} doubles the protected dimension and retains its real squared audit. Its unique optimal probability kernel must reproduce $u$. Compactness of each fixed-memory controller class and the uniform geometric word tail give attainment; equality at a finite size would contradict the exact-memory theorem. For the approximation, store $u_0,\ldots,u_{T-1}$ and thereafter output probability zero. Since $u_n\le\lambda^n$, summing the exact risk identity \eqref{eq:r67-generating-risk} yields \eqref{eq:r70-approx}. Nothing in this construction observes leakage or changes first-herald stopping.
\end{proof}

This is not the earlier unique-sign counterexample under depolarization. Here the required response is probability valued and its absolute nonconstant eigenvalues all lie strictly inside the unit disk. The impossibility comes from their phases relative to the visible decay rate. The stationary state is on a face, so the interior hypothesis of the instrument-wide realization theorem is genuinely unavailable. With $\gamma=1$ and $\varphi=\pi$, moreover, $u_0=0$ and $0\le u_n\le\lambda$ for all $n\ge1$. The one-label zero response is uniformly within $\lambda$ of this law while exact memory remains infinite for every positive $\lambda$. Exact realization is not a robust approximate-memory assertion.

\subsection{Finitely many physical models detect the obstruction}
The continuum of herald-flag parameters used to force an arbitrary full response can be replaced, against a prescribed finite-memory bound, by a finite set whose cardinality depends only on that bound and the target's linear order.

\begin{theorem}[Finite rational-resolvent tests and fixed-dimensional witnesses]\label{thm:r70-finite-tests}
Suppose a target channel response has generating function $U(z)=P(z)/D(z)$ with
$\deg P\le r-1$, $\deg D\le r$, and $D(z)\ne0$ on $[0,q\bar t]$. Fix $M\ge1$ and any $M+r$ distinct parameters $t_j\in[0,\bar t]$. In the squared embedding, an observer with at most $M$ live labels has risk exactly $2$ on the finite family $(t_j,\sigma)$, $\sigma\in\{-1,1\}$, if and only if it realizes the entire target response at every failure count.

Consequently for the algebraic leaking qubit with $\lambda=1/2$, $\alpha=\gamma=1$, $e^{i\vartheta}=(3+4i)/5$ and $\varphi=0$, the following holds for every $M$: a six-dimensional prepared squared experiment with only $2(M+3)$ rationally specified models, one unit-cost command and the fixed two-value dictionary has boundary value $2$, but every at-most-$M$-label observer has risk at least $2+\delta_M$, where $\delta_M>0$ is an attained, computable real algebraic number. One may take $p=1/2$, $\bar t=1/2$ and $t_j=j/[2(M+2)]$, $0\le j\le M+2$. Neither the physical dimension nor the target response order increases with $M$.
\end{theorem}
\begin{proof}
Boundary feasibility still forces first-herald commitment in this finite model family, since its physical herald probability is always $p$. An $M$-label rule has
$B(z)=e_1^{\mathsf T}(I-zQ)^{-1}d=P_Q(z)/D_Q(z)$,
where $\deg P_Q\le M-1$, $\deg D_Q\le M$ and $D_Q(z)\ne0$ for $0\le z<1$. Smaller matrices may be padded by unreachable rows.

The two signs and \eqref{eq:r67-generating-risk} make risk $2$ equivalent to $B(qt_j)=U(qt_j)$ at every chosen parameter. The polynomial
$P_QD-PD_Q$ has degree at most $M+r-1$ and vanishes at $M+r$ distinct points. It is therefore zero. Equality of the power series about zero gives every response coefficient. The converse follows immediately from the same risk identity. This is rational interpolation with known degree, not extrapolation from arbitrary fitted samples.

For the stated target,
\[
 U(z)=\frac1{2(1-\lambda z)}+
       \frac{1-\lambda cz}{2(1-2\lambda cz+\lambda^2z^2)},
 \qquad c=3/5,
\]
so $r=3$ suffices and all coefficients are rational. The exact leakage theorem excludes every finite realization. Hence the finite tests exclude equality for the specified $M$. The full-history target kernel still attains $2$ in every selected model. Compact controller simplexes and the geometric tail imply the fixed-$M$ minimum is attained, so its gap $\delta_M$ is strictly positive.

Finally introduce real variables $Q,d$ and vectors $v_j$ satisfying
$(I-qt_jQ)v_j=d$. These matrices are invertible for every stochastic $Q$, and the achieved risk is
$2+2\varepsilon p\max_j|(v_j)_1-U(qt_j)|$.
All defining constraints are polynomial over the rationals. Real-algebraic elimination computes the attained minimum and an algebraic minimizing controller; in particular the positive gap is algebraic and computable. This is an exact finite procedure, not a polynomial-time claim or a numerical certificate inferred from local optimization. The general procedure is proved here; a full elimination at every $M$ has not been executed in the finite diagnostics.
\end{proof}

Thus finite linear rank four does not bound the classical memory forced by a finite family of actual squared-loss tests, even on a fixed protected dimension. The number of models in this conclusion grows with the tested memory bound. It is not a claim that one fixed finite model list excludes every finite observer.

\subsection{Positive background regularizes the output face}
The leakage obstruction is sensitive to the face on which the terminal effect vanishes. Replace the effect by
$E_\eta=\eta I/2+(1-\eta)E$, $0<\eta<1$, on the same qutrit. The response becomes
\[
 u_n^{(\eta)}=\eta/2+(1-\eta)u_n.
\]
This changes a real terminal measurement, not the protected channel or the observer's information.

\begin{theorem}[Exact finite realization after background regularization]\label{thm:r70-background}
For all parameters of Theorem~\ref{thm:r70-leakage} and every $0<\eta<1$, the response $u^{(\eta)}$ has a finite exact stochastic realization. If $T\ge1$ satisfies
\begin{equation}\label{eq:r70-background-depth}
 (1-\eta)\lambda^T\le\eta/2,
\end{equation}
then at most $T+2+m(\lambda\alpha,\vartheta)$ live labels suffice. The construction uses a prescribed initial label and attains the squared-task optimum at the same cost $1/p$.

For fixed $\lambda\in(0,1)$, $\alpha=1$ and irrational phase, the least exact size tends to infinity as $\eta\downarrow0$, while the displayed upper is $O_\lambda(\log(1/\eta))$. No matching lower rate in $\eta$ is asserted. For algebraic nondegenerate inputs this finite upper and the inherited fixed-size response test give a terminating search for a minimal exact controller.
\end{theorem}
\begin{proof}
Put $a=(1-\eta)\lambda^T/2$ and $\psi=T\vartheta+\varphi$. The suffix starting at $T$ is the average of the two probability responses
\[
 f_k=\eta/2+2a\lambda^k,\qquad
 g_k=\eta/2+2a\gamma\alpha^T(\lambda\alpha)^k\cos(k\vartheta+\psi).
\]
The first has a two-label transient/absorbing realization with decoder values $\eta/2+2a$ and $\eta/2$. For the second, use the normalized eigenvector construction of Theorem~\ref{thm:r69-spectral} at eigenvalue $\lambda\alpha e^{i\vartheta}$, with decoder $\eta/2+2a\gamma\alpha^T\operatorname{Re}(e^{i\psi}v_j)$. Condition \eqref{eq:r70-background-depth} gives $2a\le\eta/2$, so all these decoder values lie in $[0,\eta]\subset[0,1]$.

Store a deterministic chain for $u_0^{(\eta)},\ldots,u_{T-1}^{(\eta)}$. At its final failure transition enter the two disjoint suffix machines with weights $1/2,1/2$. No additional initialization label is needed: the last prefix row makes and stores the choice within the counted blocks. At every later failure count the output is $(f_k+g_k)/2=u_{T+k}^{(\eta)}$. This proves the exact count and preserves the physical stopping rule.

If exact sizes stayed bounded along $\eta\downarrow0$, pad their stochastic matrices and decoders to one size and take a convergent subsequence in the compact row simplexes. Every finite coefficient is continuous in those variables. The limit would realize $u$ at $\eta=0$, contradicting Theorem~\ref{thm:r70-leakage}. The logarithmic upper follows by solving \eqref{eq:r70-background-depth} with fixed $\lambda$. For algebraic input, the target has a supplied linear representation of dimension at most four. The finite word-equality criterion of Proposition~\ref{prop:finite-realization-test}, together with this size upper, terminates a minimal exact-size search. Algebraic controller encodings and exact sampling resources remain distinct from label count.
\end{proof}

\section{The causal task-order transfer theorem}\label{sec:central}
Suppose that every terminating state-stopped competitor reaches at least $B_0\ge1$ live labels under $\theta_*$. Such a certificate must concern actual physical command compatibility, not printed action names. Define
\begin{equation}\label{eq:widths}
\begin{array}{c|cc}
 i&m_-(M)&m_+(M)\\\hline
 \cp&M&\lfloor M/D_0\rfloor\\
 \ev&M&\lfloor M/D_0\rfloor\\
 \st&M-B_0&\lfloor(M-B)/D_0\rfloor .
\end{array}
\end{equation}
The event upper also requires $M\ge B$. The trivial state-stopped lower $B_0=1$ is valid for a procedure initially waiting to make a call.

\begin{theorem}[Affine task-order transfer with a paid reference and common construction]\label{thm:transfer}
Assume all the prepared-kernel, full-report, task-order, effort, occupation, resource and common-construction hypotheses of Section~\ref{sec:data}. Fix $N,M$, put $b_-=N/\kappa_-$, $b_+=N/\kappa_+$, and suppose $b_-\ge1$ and $m_-(M)\ge1$. Set
\[
 t_N=(1-A_N)_+/u_-.
\]
If $t_N>0$, every one of the three risks satisfies
\begin{equation}\label{eq:central-lower}
 R_i^\Theta(N,M)\ge
 \bigl[l_-t_N\Q_{m_-(M)}(b_-/t_N)-E_N\bigr]_+.
\end{equation}
If $t_N=0$, the asserted lower is zero, with no evaluation of $\Q$ at infinity. An empty legal class still has risk $+\infty$.

Suppose also $m_+(M)\ge1$, $b_+\ge1$, and, for event stopping, $M\ge B$. If
\begin{equation}\label{eq:beta}
 \beta_N=\frac{l_+b_+}{1+b_+a_+}\ge1,
\end{equation}
one common Borel program with actual expected cost at most $N$ in every model attains
\begin{equation}\label{eq:central-upper}
 R_i^\Theta(N,M)\le
 \frac{u_+(1+b_+a_+)}{l_+}\Q_{m_+(M)}(\beta_N)+b_+e_+.
\end{equation}
It commits almost surely. If $\beta_N<1$, the always-accept first-attempt construction instead gives $R_i^\Theta(N,M)\le u_+\Q_{m_+(M)}(1)+e_+$ at cost at most $\kappa_+\le N$.

The lower is against every legal competing strategy and uses one actual witness. The upper is one parameter-independent strategy for all models. The two effort constants need not agree. Neither direction assumes that the comparison experiment belongs to $\Theta$ or that a normalized predictive update is uniformly contracting.
\end{theorem}

\begin{corollary}[Scale-uniform relative transfer]\label{cor:relative}
Under Theorem~\ref{thm:transfer}, suppose $A_N=E_N=a_+=e_+=0$. If $l_+b_+\ge1$, then
\begin{equation}\label{eq:relative}
 \frac{l_-}{u_-}\Q_{m_-(M)}(u_-b_-)
 \le R_i^\Theta(N,M)
 \le\frac{u_+}{l_+}\Q_{m_+(M)}(l_+b_+).
\end{equation}
The factors $l_\pm,u_\pm$ may be fixed nonunit constants, independent of budget and resolution. In particular no term proportional to $b_\pm$ times a fixed absolute calibration error is required in this relative-order class.
\end{corollary}
\begin{proof}
Substitute $t_N=1/u_-$ and $\beta_N=l_+b_+$ in the theorem.
\end{proof}

\begin{corollary}[Separate exact checkpoint and autonomous laws]\label{cor:exact}
Suppose the certificates are exact: $l_\pm=u_\pm=1$, all defects vanish, $\kappa_-=\kappa_+=\kappa$, and exports are the identity. Then, for $N\ge\kappa$,
\[
 R_\cp^\Theta(N,m)=\Q_m(N/\kappa).
\]
The constructed autonomous upper uses $B+m$ state-stopped labels, or $\max(B,m)$ event-stopped labels. If the mandatory live count is sharp, $B_0=B$, then
\begin{align}
 R_\st^\Theta(N,B+m)&=\Q_m(N/\kappa)\quad(m\ge1),\label{eq:typed-exact}\\
 R_\ev^\Theta(N,M)&=\Q_M(N/\kappa)\quad(M\ge B).\label{eq:event-exact}
\end{align}
The subtraction by $B$ belongs only to the state-stopped model. Nonidentity exports require the expanded dictionary in \eqref{eq:widths}, unless an optimal dictionary is invariant under all exports. These exact stopping distinctions are inherited from the complete companion and are reproved here to fix the resource semantics.
\end{corollary}
\begin{proof}
The same lower and upper profile coincide in each asserted interface. Identity exports give $D_0=1$. The packing and converse widths are proved in Lemma~\ref{lem:roles}. No lower on live states is used to forbid event-stopped reuse.
\end{proof}

The familiar absolute-defect theorem is also a specialization: take $l_\pm=u_\pm=1$ while retaining $A_N,E_N,a_+,e_+$. Thus relative and absolute calibration are composed in one statement. The chosen metric is the common scored task throughout; it is not changed between the converse and the construction.

\section{Proof of the central theorem}
\begin{lemma}[Profile attainment]\label{lem:attainment}
For every finite $m\ge1$ and $b\ge1$, the minimum in \eqref{eq:profile} is attained. The profile is nonincreasing in $m,b$ and left-continuous in $b>1$.
\end{lemma}
\begin{proof}
Write a dictionary as an ordered tuple in the compact set $H^m$, permitting repetitions. The feasible densities are a weak-star closed subset of the radius-$b$ ball of $L^\infty(\nu)$: positivity, the upper bound and the mass equality are each closed under pairing with $L^1(\nu)$. They are therefore weak-star compact. As the tuple varies, its squared-distance function varies uniformly on $H$, since all points belong to a bounded set. Along a convergent product net, split the objective difference into this uniform variation and the weak-star pairing against the fixed bounded measurable loss at the limit tuple. This proves continuity of the objective and attainment. Nets avoid an unwarranted sequential compactness assumption. Monotonicity follows from inclusion of feasible sets. For left continuity take an optimizer $h$ at $b>1$ and let $h_\lambda=(1-\lambda)h+\lambda$. It has unit mass and cap $b_\lambda=(1-\lambda)b+\lambda<b$. Its objective tends to the optimum as $\lambda\downarrow0$, while monotonicity bounds the optimum at $b_\lambda$ between the optimum at $b$ and this objective.
\end{proof}

\begin{lemma}[Actual-history comparison measure]\label{lem:comparison}
Fix any robustly admissible competitor of cost at most $N$, with its fixed terminal dictionary $C$. There are finite nonnegative measures $\xi_c$ on $X$, $c\in C$, such that $\xi=\sum_c\xi_c$ obeys
\begin{align}
 \xi&\le (\E_*A_\tau)\nu\le b_-\nu,\label{eq:measure-cap}\\
 1&\le u_-\xi(X)+A_N,\label{eq:measure-mass}\\
 \E_*\|Z-\widehat r\|^2
 &\ge l_-\sum_{c\in C}\int d_c\,d\xi_c-E_N.
 \label{eq:measure-risk}
\end{align}
The measure $\xi$ is not asserted to be the terminal law of any comparison execution.
\end{lemma}
\begin{proof}
Adjoin an independent uniform random variable to the current report on both sides. This is current-call randomness, not a seed retained between calls; it is sufficient to represent the competitor's randomized output and stopping decision. A finite-output kernel has such a representation by cumulative probabilities. At a fixed reached productive history $h$ and chosen command $a$, the post-report decision then gives an acceptance indicator $q_h(y)$ and an output $r_h(y)\in C$ on acceptance. Give $r_h$ an arbitrary fixed value on rejection. The notation $y$ now includes the current independent random variable. Integrating the original certificates over that variable preserves their constants. Dependence on the entire earlier history is permitted in these functions.

For each Borel $D\subset X$, set
\begin{equation}\label{eq:xi-definition}
 \xi_c(D)=\E_*\sum_{j\le\tau:\ j\ {
m productive}}
 \int\one_{\{\pi_{H_{j-1},a_j}(y)\in D\}}
 q_{H_{j-1}}(y)\one_{\{r_{H_{j-1}}(y)=c\}}
 P^0_{H_{j-1},a_j}(dy,H).
\end{equation}
The outer expectation is the actual pre-call occupation at $\theta_*$. All terms are jointly measurable and nonnegative. Tonelli defines the measures, whose total is finite because $\E_*A_\tau\le b_-$. Summing over $c$ and dropping the acceptance indicator gives \eqref{eq:measure-cap}: each comparison row has mark law $\nu$.

The analogous sum with the actual rows instead of $P^0$ has mass exactly one, because the actual procedure commits exactly once almost surely. Apply \eqref{eq:lower-mass} at each actual reached history, then sum. The occupation certificate yields \eqref{eq:measure-mass}. Apply \eqref{eq:lower-task} in the same way to the disjoint actual commitment events. Their expected losses sum to the actual risk. The comparison contribution is $d_{r_h(y)}(\pi_{h,a}y)$ by \eqref{eq:full-moments}; this identity conditions on the full current report and therefore does not overlook extra information seen by the competitor. The accumulated loss defect is at most $E_N$, giving \eqref{eq:measure-risk}. One may first truncate the number of calls. Monotone convergence for the nonnegative contributions and the finite defect bounds justify the limit.

No actual past is replaced by a comparison past. In particular the competitor need not terminate, or obey a cost bound, if hypothetically run under the comparison experiment. This distinction is why the argument applies when that experiment is not an operating model.
\end{proof}

\begin{remark}[Why the full report is required]
A constant mark can conceal a visible fair bit which also equals the future audit. A comparison using only that mark has $f=0,w=1$, but a two-value decoder reading the bit predicts it with zero loss. The hypotheses therefore condition on the full report seen by the competitor; matching a compressed mark marginal alone cannot establish the lower.
\end{remark}

\begin{lemma}[Live roles, output roles and physical compatibility]\label{lem:roles}
A program with $\ell$ live roles and $t$ output roles has a state-stopped implementation with $\ell+t$ labels and an event-stopped implementation with $\max(\ell,t)$ labels, preserving every path cost and complete stopped scored law. Conversely an event-stopped $M$-label program has at most $M$ decoder values. A state-stopped $M$-label program with at least $B_0$ reached live labels has at most $M-B_0$ terminal decoder values.

For a physical live-state certificate, consider finitely many mandatory pre-call contexts $I$ with no free context input. Suppose the commands legal at a history of context $i$ are exactly $\mathcal U_i$, and every terminating competitor visits every context with positive probability under the same witness. This positive probability may depend on the competitor; no uniform lower bound over competitors is required. Let $b_{\rm phys}$ be the least size of a cover of $I$ by subsets $J$ with $\bigcap_{i\in J}\mathcal U_i\ne\varnothing$. Then at least $b_{\rm phys}$ live labels are required. This certificate is invariant under duplicate names and bijective stateless recodings of the same physical commands.
\end{lemma}
\begin{proof}
For state stopping take disjoint copies of the live and terminal role sets and type them zero and one. For event stopping embed the live and terminal role sets separately into one alphabet of size $\max(\ell,t)$. On continuing transitions use the live embedding; on commitment emit the final event and use the terminal embedding. A command row is read only on a continuing branch, while the decoder is read only after commitment. Sharing their label therefore causes no conflict and uses no continuing mode bit. Couple the original and recoded coins. Every finite stopped prefix and every path cost agree, and almost-sure termination gives equality of the stopped scored laws. A fixed decoder has one value per retained label. In state stopping the global halt map excludes each reached live label from the output set, proving its subtraction. Event stopping has no such exclusion.

A live label reached in contexts $J$ must have one command distribution supported on all the corresponding $\mathcal U_i$. Since $J$ is finite, the union of the exceptional null command sets is null; hence their intersection supports a probability and is nonempty. The sets of contexts served by the labels cover $I$, proving the bound. Conversely a nonempty intersection contains a command and hence a Dirac command distribution. Replacing names without changing executed kernels, legal domains, costs and reports does not change these intersections. A translator which itself remembers a hidden phase is not such a recoding; its memory must be charged.
\end{proof}

\begin{proof}[Proof of Theorem~\ref{thm:transfer}]
Fix a competing program. Lemma~\ref{lem:roles} bounds its decoder width by $m_-(M)$. Let $s=\xi(X)$ in Lemma~\ref{lem:comparison}. If $t_N>0$, then \eqref{eq:measure-mass} gives $s\ge t_N$. Scale the measure to $\xi'=(t_N/s)\xi$. It has mass $t_N$ and remains bounded above by $b_-\nu$. Thus $h=(d\xi'/d\nu)/t_N$ has unit mass and cap $b_-/t_N$. Nearest-dictionary projection and nonnegativity imply
\[
 \sum_{c\in C}\int d_c\,d\xi_c
 \ge\int [w+\dist(f,C)^2]\,d\xi
 \ge t_N\Q_{m_-(M)}(b_-/t_N).
\]
Together with \eqref{eq:measure-risk} this proves the asserted lower for this competitor. Its worst-model risk dominates its risk at $\theta_*$. Taking the infimum over competitors gives \eqref{eq:central-lower}, without exchanging an infimum and a supremum. Nonnegative loss gives the zero truncation and the stated case $t_N=0$.

For the upper assume \eqref{eq:beta}. Choose an attained optimizer $(C,h)$ at width $m_+(M)$ and cap $\beta_N$. Let $r_C(x)$ be the first nearest point to $f(x)$ in a fixed ordering of $C$. After each attempt of $\mathcal G$, accept $x$ with probability $h(x)/\beta_N$, export $T_Jr_C(x)$ and commit; otherwise return to the start routing label. All these operations are Borel. The choice $J$ is used within the current attempt; any part retained thereafter is charged through the live role or exported decoder, not a free seed. The exports are assumed jointly measurable and take at most $D_0|C|$ distinct values, so their finite decoder and transitions are Borel as well.

Conditional on every surviving attempt history, the acceptance probability is at least
\begin{equation}\label{eq:hazard}
 p=\frac{l_+}{\beta_N}-a_+=\frac1{b_+}>0.
\end{equation}
Consequently the probability of reaching attempt $j+1$ is at most $(1-p)^j$, by induction on conditional survival probabilities. In particular commitment is almost sure and the expected number of attempts is at most $1/p=b_+$. Each begun attempt has conditional expected cost at most $\kappa_+$, not merely a cost conditional on a favorable outcome. Summing the predictable attempt costs gives actual expected cost at most $\kappa_+b_+=N$.

The conditional unnormalized loss contributed by an attempt is at most $u_+\Q_{m_+(M)}(\beta_N)/\beta_N+e_+$. Sum its contributions on the disjoint terminal events, using the same survival bound. The result is at most
\[
 b_+\left[\frac{u_+}{\beta_N}\Q_{m_+(M)}(\beta_N)+e_+\right],
\]
which is \eqref{eq:central-upper}. The program has $B$ live roles and at most $D_0m_+(M)$ terminal roles. Lemma~\ref{lem:roles} gives exactly the sufficient widths in \eqref{eq:widths}. A checkpoint counts only the terminal width. No physical renewal law was used: the conditional hazard controls survival even if the hidden state after rejection differs arbitrarily from its former state.

If $\beta_N<1$, commit after the first attempt, and use an optimizer at cap one. The loss certificate with $q=1$ gives $u_+\Q_{m_+(M)}(1)+e_+$, and the whole-attempt cost is at most $\kappa_+$. This proves all remaining assertions.
\end{proof}

\section{Scope and continuing mathematical questions}\label{sec:r71-scope}
The arithmetic theorem determines a constant-factor exact-memory envelope for each nonreal single-frequency probability response, and both logarithmic growth exponents for every irrational phase. Its uniform cubic-root criterion is necessary and sufficient. It does not determine the irrationality exponent of the particular angle with $e^{i\vartheta}=(3+4i)/5$, nor a sharp leading memory constant. In a finite mixture it is the maximum of the component envelopes that governs memory; simultaneous favorable approximation scales are not assumed.

The exact realization uses a prescribed initial live label, stochastic transitions and a bounded decoder. Its physical squared-task interpretation retains the original first-herald cost $1/p$, fixed dictionary and actual quantum audit. Transition-probability descriptions, exact sampling, program storage and arithmetic time are additional resources. An ideal exact-label theorem is not a matching complexity theorem or an algorithm for a noncomputable phase.

The current supplement contains the full instrument-wide state-interior realization, normalized-budget stopping rigidity, robust tied optimality and finite squared-witness proofs. They remain independent results with their own scopes: a latent instrument simulator is not a free observed state, and the boundary criterion is not a classification of several-command feedback control. The inherited performed-law, minimax, reset, timing, finite-width, rank and geometric consequences are retained in full. The separate nonintervening sensing, ordered-measurement, two-stage reset communication and rank-constrained discrimination lines remain at their immutable source identities.

The geometric transfer still requires its own actual-history lower witness and one common paid construction. General unknown-model predictive realization, arbitrary controlled boundary stopping, the sharp general-$M$ squared-sign and phase-type timing risk frontiers, and matched program, precision, workspace and time lower bounds remain distinct research questions. The original title and top-level research objective are retained without identifying the present phase classification with completion of the entire programme.
\begin{thebibliography}{99}
\bibitem{Blackwell} D. Blackwell, \emph{Equivalent comparisons of experiments}, Ann. Math. Statist. \textbf{24} (1953), 265--272. DOI: 10.1214/aoms/1177729032.
\bibitem{CGM97} J. A. Cuesta-Albertos, A. Gordaliza and C. Matr\'an, \emph{Trimmed $k$-means: an attempt to robustify quantizers}, Ann. Statist. \textbf{25} (1997), 553--576. DOI: 10.1214/aos/1031833664.
\bibitem{GG99} L. A. Garc\'ia-Escudero and A. Gordaliza, \emph{Robustness properties of $k$ means and trimmed $k$ means}, J. Amer. Statist. Assoc. \textbf{94} (1999), 956--969. DOI: 10.1080/01621459.1999.10474200.
\bibitem{SC} C. R. Shalizi and J. P. Crutchfield, \emph{Computational mechanics: pattern and prediction, structure and simplicity}, J. Statist. Phys. \textbf{104} (2001), 817--879. DOI: 10.1023/A:1010388907793.
\bibitem{BC} N. Barnett and J. P. Crutchfield, \emph{Computational mechanics of input--output processes: structured transformations and the $\varepsilon$-transducer}, J. Statist. Phys. \textbf{161} (2015), 404--451; arXiv:1412.2690.
\bibitem{AMTW} A. Ambainis, M. Mosca, A. Tapp and R. de Wolf, \emph{Private quantum channels}, Proc. 41st Annual Symposium on Foundations of Computer Science (2000), 547--553. DOI: 10.1109/SFCS.2000.892142.
\bibitem{Buscemi} F. Buscemi, \emph{Comparison of quantum statistical models: equivalent conditions for sufficiency}, Comm. Math. Phys. \textbf{310} (2012), 625--647. DOI: 10.1007/s00220-012-1421-3.
\bibitem{BR14} N. Bertschinger and J. Rauh, \emph{The Blackwell relation defines no lattice}, Proc. IEEE International Symposium on Information Theory (2014), 2479--2483. DOI: 10.1109/ISIT.2014.6875280; arXiv:1401.3146.
\bibitem{Banerjee05} A. Banerjee, X. Guo and H. Wang, \emph{On the optimality of conditional expectation as a Bregman predictor}, IEEE Trans. Inform. Theory \textbf{51} (2005), 2664--2669. DOI: 10.1109/TIT.2005.850145.
\bibitem{Stengel96} B. von Stengel, \emph{Efficient computation of behavior strategies}, Games Econom. Behav. \textbf{14} (1996), 220--246. DOI: 10.1006/game.1996.0050.
\bibitem{Balder89} E. J. Balder, \emph{On compactness of the space of policies in stochastic dynamic programming}, Stochastic Process. Appl. \textbf{32} (1989), 141--150. DOI: 10.1016/0304-4149(89)90058-6.
\bibitem{Barry14} J. Barry, D. T. Barry and S. Aaronson, \emph{Quantum partially observable Markov decision processes}, Phys. Rev. A \textbf{90} (2014), 032311. DOI: 10.1103/PhysRevA.90.032311; arXiv:1406.2858.
\bibitem{Ying18} S. Ying and M. Ying, \emph{Reachability analysis of quantum Markov decision processes}, Inform. Comput. \textbf{263} (2018), 31--51. DOI: 10.1016/j.ic.2018.09.001; arXiv:1406.6146.
\bibitem{Saldi25} N. Saldi, S. Sanjari and S. Y\"uksel, \emph{Quantum Markov decision processes: general theory, approximations, and classes of policies}, SIAM J. Control Optim. \textbf{63} (2025), 1156--1183. DOI: 10.1137/24M1643864; arXiv:2402.14649.
\bibitem{Seidenberg54} A. Seidenberg, \emph{A new decision method for elementary algebra}, Ann. of Math. (2) \textbf{60} (1954), 365--374. DOI: 10.2307/1969640.
\bibitem{Chatterjee15} K. Chatterjee, M. Chmel\'ik, R. Gupta and A. Kanodia, \emph{Optimal cost almost-sure reachability in POMDPs}, Proc. AAAI Conference on Artificial Intelligence \textbf{29} (2015). DOI: 10.1609/aaai.v29i1.9683.
\bibitem{Andriushchenko23} R. Andriushchenko, A. Bork, M. \v{C}e\v{s}ka, S. Junges, J.-P. Katoen and M. Mac\'ak, \emph{Search and explore: symbiotic policy synthesis in POMDPs}, Computer Aided Verification (2023); arXiv:2305.14149v2.
\bibitem{Halmos69} P. R. Halmos, \emph{Two subspaces}, Trans. Amer. Math. Soc. \textbf{144} (1969), 381--389. DOI: 10.1090/S0002-9947-1969-0251519-5.
\bibitem{Kessler21} D. A. Kessler, E. Barkai and K. Ziegler, \emph{First-detection time of a quantum state under random probing}, Phys. Rev. A \textbf{103} (2021), 022222. DOI: 10.1103/PhysRevA.103.022222; arXiv:2012.01763.
\bibitem{DelVecchio26} G. Del Vecchio Del Vecchio and S. N. Majumdar, \emph{Optimal detection of quantum states via projective measurements}, J. Phys. A: Math. Theor. \textbf{59} (2026), 035001. DOI: 10.1088/1751-8121/ae34fe; arXiv:2509.08556.
\bibitem{Companion} Q. Qi, \emph{General Theta Foundations I: Acquired Geometry and Causal Resource Transfer}, complete technical companion on absolute calibration, exact stopping resources and nonintervening sensing, 2026. The complete retained R58 native source and its independently built companion accompany this manuscript; their identities are fixed in the source-provenance and preservation records.
\bibitem{Tzeng92} W.-G. Tzeng, \emph{A polynomial-time algorithm for the equivalence of probabilistic automata}, SIAM J. Comput. \textbf{21} (1992), 216--227. DOI: 10.1137/0221017.
\bibitem{LiYuYing12} Y. Li, N. Yu and M. Ying, \emph{Termination of nondeterministic quantum programs}, arXiv:1201.0891 (2012).
\bibitem{Fu24} J. Fu, H. Jiang, M. Xu, Y. Deng and Z.-B. Li, \emph{Algorithmic analysis of termination problems for nondeterministic quantum programs}, arXiv:2402.15827 (2024).

\bibitem{LiuVerdu18} J. Liu and S. Verd\'u, \emph{Rejection sampling and noncausal sampling under moment constraints}, Proc. IEEE International Symposium on Information Theory (2018), 1565--1569. DOI: 10.1109/ISIT.2018.8437857.
\bibitem{GocFlamich24} D. Goc and G. Flamich, \emph{On channel simulation with causal rejection samplers}, Proc. IEEE International Symposium on Information Theory (2024). DOI: 10.1109/ISIT57864.2024.10619339; examined text arXiv:2401.16579v2 (3 May 2024).
\bibitem{LiLiu26} Z. Li and T. Liu, \emph{Ranking experiments under sequential sampling}, arXiv:2608.19897v2 (28 August 2026). Version 1 appeared under the title \emph{The order of binary experiments under endogenous stopping}, with Z. Li as sole author; the version-2 statement is the one compared here.
\bibitem{Ying14pre} S. Ying and M. Ying, \emph{Reachability analysis of quantum Markov decision processes}, arXiv:1406.6146 (2014), Theorem 3.6; published journal version listed above as \cite{Ying18}.

\bibitem{KaraYuksel20} A. D. Kara and S. Y\"uksel, \emph{Robustness to incorrect system models in stochastic control}, SIAM J. Control Optim. \textbf{58} (2020), 1144--1182. DOI: 10.1137/18M1208058; examined text arXiv:1803.06046v7.
\bibitem{Demirci25} Y. E. Demirci, A. D. Kara and S. Y\"uksel, \emph{Sensitivity of filter kernels and robustness bounds to transition and measurement kernel perturbations in partially observable stochastic control}, arXiv:2508.10658v3 (2025).
\bibitem{Watrous18} J. Watrous, \emph{The Theory of Quantum Information}, Cambridge University Press, 2018, Chapter 3, Corollary 3.40.
\bibitem{KRN65} K. Kuratowski and C. Ryll-Nardzewski, \emph{A general theorem on selectors}, Bull. Acad. Polon. Sci. S\'er. Sci. Math. Astronom. Phys. \textbf{13} (1965), 397--403.

\bibitem{Charnes63} A. Charnes, W. W. Cooper and K. O. Kortanek, \emph{Duality in semi-infinite programs and some works of Haar and Carath\'eodory}, Management Sci. \textbf{9} (1963), 209--228. DOI: 10.1287/mnsc.9.2.209.
\bibitem{Shapiro09} A. Shapiro, \emph{Semi-infinite programming, duality, discretization and optimality conditions}, Optimization \textbf{58} (2009), 133--161. DOI: 10.1080/02331930902730070; author preprint, Optimization Online (2008).
\bibitem{Dufour25} F. Dufour and T. Prieto-Rumeau, \emph{Absorbing Markov decision processes and their occupation measures}, SIAM J. Control Optim. \textbf{63} (2025), 676--698. DOI: 10.1137/24M1661959. Related author preprint examined: \emph{A note on weak compactness of occupation measures for an absorbing Markov decision process}, arXiv:2402.10672v2 (2024).
\bibitem{Franc23} V. Franc, D. Prusa and V. Voracek, \emph{Optimal strategies for reject option classifiers}, J. Mach. Learn. Res. \textbf{24} (2023), no. 11, 1--49.
\bibitem{Jencova16} A. Jen\v{c}ov\'a, \emph{Comparison of quantum channels and statistical experiments}, arXiv:1512.07016v2 (2016). The separately indexed shorter version arXiv:1601.06370 was withdrawn in favor of this extended text.
\bibitem{LiElGamal18} C. T. Li and A. El Gamal, \emph{Strong functional representation lemma and applications to coding theorems}, IEEE Trans. Inform. Theory \textbf{64} (2018), 6967--6978; arXiv:1701.02827v4.
\bibitem{AldousShepp87} D. Aldous and L. Shepp, \emph{The least variable phase type distribution is Erlang}, Comm. Statist. Stochastic Models \textbf{3} (1987), 467--473. DOI: 10.1080/15326348708807067.
\bibitem{Sion58} M. Sion, \emph{On general minimax theorems}, Pacific J. Math. \textbf{8} (1958), 171--176. DOI: 10.2140/pjm.1958.8.171.
\bibitem{PiunovskiyZhang24} A. Piunovskiy and Y. Zhang, \emph{Extreme occupation measures in Markov decision processes with an absorbing state}, SIAM J. Control Optim. \textbf{62} (2024), 65--90. DOI: 10.1137/23M1572398. Author text examined: \emph{Extreme occupation measures in Markov decision processes with a cemetery}, arXiv:2307.03158v1.
\bibitem{DufourMixtures} F. Dufour and T. Prieto-Rumeau, \emph{Absorbing Markov decision processes: geometric properties and sufficiency of finite mixtures of deterministic policies}, arXiv:2512.17511; examined HTML at version-1 URL, Theorems 4.3--4.4 and Proposition 4.5.
\bibitem{Mereghetti24} C. Mereghetti, B. Palano and P. Raucci, \emph{Unary quantum finite state automata with control language}, Appl. Sci. \textbf{14} (2024), no. 4, 1490. DOI: 10.3390/app14041490.
\bibitem{RabinScott59} M. O. Rabin and D. Scott, \emph{Finite automata and their decision problems}, IBM J. Res. Develop. \textbf{3} (1959), 114--125. DOI: 10.1147/rd.32.0114.
\bibitem{Jaeger00} H. Jaeger, \emph{Observable operator models for discrete stochastic time series}, Neural Comput. \textbf{12} (2000), 1371--1398. DOI: 10.1162/089976600300015411.
\bibitem{ThonJaeger15} M. Thon and H. Jaeger, \emph{Links between multiplicity automata, observable operator models and predictive state representations---a unified learning framework}, J. Mach. Learn. Res. \textbf{16} (2015), 103--147.
\bibitem{MonrasWinter14} A. Monras and A. Winter, \emph{Quantum learning of classical stochastic processes: the completely-positive realization problem}, LIPIcs \textbf{27}, TQC 2014 (2014), 99--109. DOI: 10.4230/LIPIcs.TQC.2014.99.

\bibitem{Heller65} A. Heller, \emph{On stochastic processes derived from Markov chains}, Ann. Math. Statist. \textbf{36} (1965), 1286--1291. DOI: 10.1214/aoms/1177700000.
\bibitem{CarlylePaz71} J. W. Carlyle and A. Paz, \emph{Realizations by stochastic finite automata}, J. Comput. System Sci. \textbf{5} (1971), 26--40. DOI: 10.1016/S0022-0000(71)80005-3.
\bibitem{Fanizza24} M. Fanizza, J. Lumbreras and A. Winter, \emph{Quantum theory in finite dimension cannot explain every general process with finite memory}, Comm. Math. Phys. \textbf{405} (2024), article 50. DOI: 10.1007/s00220-023-04913-4; arXiv:2209.11225. In particular see Theorem 2 and Section 4.2 for the noise-dependent worst-phase classical-memory lower bound.
\bibitem{DmitrievDynkin46} N. Dmitriev and E. Dynkin, \emph{On characteristic roots of stochastic matrices}, Izv. Akad. Nauk SSSR Ser. Mat. \textbf{10} (1946), no. 2, 167--184. Math-Net.Ru: im3595 (Russian).
\bibitem{JohnsonPaparella17} C. R. Johnson and P. Paparella, \emph{A matricial view of the Karpelevi\v c theorem}, Linear Algebra Appl. \textbf{520} (2017), 1--15. DOI: 10.1016/j.laa.2017.01.009; arXiv:1611.06970v2.
\bibitem{KirklandSmigoc22} S. Kirkland and H. \v Smigoc, \emph{Stochastic matrices realising the boundary of the Karpelevi\v c region}, Linear Algebra Appl. \textbf{635} (2022), 116--138. DOI: 10.1016/j.laa.2021.11.016; arXiv:2110.01040v1.
\bibitem{BenvenutiFarina04} L. Benvenuti and L. Farina, \emph{A tutorial on the positive realization problem}, IEEE Trans. Automat. Control \textbf{49} (2004), 651--664. DOI: 10.1109/TAC.2004.826715.
\bibitem{BenvenutiFarina99} L. Benvenuti and L. Farina, \emph{An example of how positivity may force realizations of ``large'' dimension}, Systems Control Lett. \textbf{36} (1999), 261--266. DOI: 10.1016/S0167-6911(98)00098-X.
\bibitem{OCinneide90} C. A. O'Cinneide, \emph{Characterization of phase-type distributions}, Comm. Statist. Stochastic Models \textbf{6} (1990), 1--57. DOI: 10.1080/15326349908807134.
\bibitem{Maier91} R. S. Maier, \emph{The algebraic construction of phase-type distributions}, Comm. Statist. Stochastic Models \textbf{7} (1991), 573--602. DOI: 10.1080/15326349108807207.
\bibitem{HoneVarona19} A. N. W. Hone and J. L. Varona, \emph{Continued fractions and irrationality exponents for modified Engel and Pierce series}, Monatsh. Math. \textbf{190} (2019), 501--516. DOI: 10.1007/s00605-018-1244-1. Author version arXiv:1810.10127v1, equation (3.3).
\end{thebibliography}

\end{document}
